\documentclass[11pt,a4paper]{article}
\usepackage[margin=2.4cm]{geometry}
\usepackage{amsmath,amssymb,amsthm,mathtools}
\usepackage[protrusion=true,expansion=false]{microtype}
\usepackage{booktabs,enumitem}
\usepackage{placeins}
\usepackage{tikz}
\usetikzlibrary{arrows.meta,positioning,calc,fit,backgrounds}
\definecolor{tierspend}{RGB}{200,60,60}
\definecolor{tierfull}{RGB}{210,140,40}
\definecolor{tierin}{RGB}{60,130,180}
\definecolor{tierrecv}{RGB}{70,150,90}
\usepackage[colorlinks=true,linkcolor=blue!50!black,citecolor=blue!50!black,urlcolor=blue!50!black]{hyperref}

\emergencystretch=3em
\hbadness=10000
\hfuzz=2pt

\newtheorem{theorem}{Theorem}[section]
\newtheorem{lemma}[theorem]{Lemma}
\newtheorem{proposition}[theorem]{Proposition}
\newtheorem{corollary}[theorem]{Corollary}
\theoremstyle{definition}
\newtheorem{definition}[theorem]{Definition}
\newtheorem{construction}[theorem]{Construction}
\newtheorem{assumption}[theorem]{Assumption}
\newtheorem{remark}[theorem]{Remark}
\newtheorem{example}[theorem]{Example}

\newcommand{\F}{\mathbb{F}}
\newcommand{\Z}{\mathbb{Z}}
\newcommand{\N}{\mathbb{N}}
\newcommand{\G}{\mathbb{G}}
\newcommand{\E}{\mathbb{E}}
\newcommand{\PP}{\mathbb{P}}
\renewcommand{\Pr}{\mathrm{Pr}}
\newcommand{\Fp}{\F_p}
\newcommand{\Fq}{\F_q}
\newcommand{\ip}[2]{\langle #1,\, #2 \rangle}
\newcommand{\Comm}{\mathsf{Commit}}
\newcommand{\Open}{\mathsf{Open}}
\newcommand{\Setup}{\mathsf{Setup}}
\newcommand{\Verify}{\mathsf{Verify}}
\newcommand{\negl}{\mathsf{negl}}
\newcommand{\poly}{\mathsf{poly}}
\newcommand{\PPT}{\mathsf{PPT}}
\newcommand{\Adv}{\mathcal{A}}
\newcommand{\Prov}{\mathcal{P}}
\newcommand{\Ver}{\mathcal{V}}
\newcommand{\Sim}{\mathcal{S}}
\newcommand{\Ext}{\mathcal{E}}
\newcommand{\Rel}{\mathcal{R}}
\newcommand{\Lang}{\mathcal{L}}
\newcommand{\nfsym}{\mathsf{nf}}
\newcommand{\cmsym}{\mathsf{cm}}
\newcommand{\cmx}{\mathsf{cmx}}
\newcommand{\cvsym}{\mathsf{cv}}
\newcommand{\rk}{\mathsf{rk}}
\newcommand{\Pallas}{\ensuremath{\mathsf{Pallas}}}
\newcommand{\Vesta}{\ensuremath{\mathsf{Vesta}}}
\newcommand{\Ep}{\ensuremath{\mathcal{E}}}
\newcommand{\NoteCommit}{\ensuremath{\mathsf{NoteCommit}}}
\newcommand{\Sinsemilla}{\ensuremath{\mathsf{Sinsemilla}}}
\newcommand{\ShortCommit}{\ensuremath{\mathsf{ShortCommit}}}
\newcommand{\CommitIvk}{\ensuremath{\mathsf{Commit}^{\mathsf{ivk}}}}
\newcommand{\Extract}{\ensuremath{\mathsf{Extract}}}
\newcommand{\Acc}{\ensuremath{\mathsf{Acc}}}
\newcommand{\GroupHash}{\ensuremath{\mathsf{GroupHash}}}
\newcommand{\ask}{\ensuremath{\mathsf{ask}}}
\newcommand{\aksym}{\ensuremath{\mathsf{ak}}}
\newcommand{\nksym}{\ensuremath{\mathsf{nk}}}
\newcommand{\ivk}{\ensuremath{\mathsf{ivk}}}
\newcommand{\fvk}{\ensuremath{\mathsf{fvk}}}
\newcommand{\ovk}{\ensuremath{\mathsf{ovk}}}
\newcommand{\dk}{\ensuremath{\mathsf{dk}}}
\newcommand{\pkd}{\ensuremath{\mathsf{pk_d}}}
\newcommand{\gd}{\ensuremath{\mathsf{g_d}}}
\newcommand{\rcv}{\ensuremath{\mathsf{rcv}}}
\newcommand{\rcm}{\ensuremath{\mathsf{rcm}}}
\newcommand{\rivk}{\ensuremath{\mathsf{r_{ivk}}}}
\newcommand{\rsk}{\ensuremath{\mathsf{rsk}}}

\renewenvironment{abstract}{%
  \small\begin{center}{\bfseries\abstractname}\end{center}\medskip\noindent\ignorespaces
}{\par\medskip}

\title{\textbf{\Huge Tachyon Guide}\\[6pt]\large The Tachyon scaling architecture, at design stage: what is specified, what is designed, what is open}
\author{m@rek.onl}
\date{}
\begin{document}
\maketitle
\thispagestyle{empty}
\begin{abstract}
This volume of \emph{The Zcash Arboretum} documents a protocol that does not yet exist.
Tachyon is the proposed rearchitecture of Zcash payments around proof-carrying
wallet state: the wallet folds each block's effects into a recursive proof, so
synchronisation ceases to scale with chain length, and the chain ceases to
carry the wallet's burdens. Because the design is in motion, this volume's
organising device is a strict classification: every claim is
\emph{specified}, \emph{designed but unspecified}, or an \emph{open problem}.
The strongest chapter is the last kind. The reader should expect this volume
to age faster than its siblings and to be revised as the design lands.
\end{abstract}
\clearpage
\tableofcontents
\clearpage
\section{Scope, sources, and method}\label{sec:tachyon-scope}

Project Tachyon is a proposed redesign of Zcash's shielded protocol around
three ideas: wallet state becomes proof-carrying data, so a wallet proves the
validity and unspentness of its own notes and a verifier checks a proof rather
than replaying history; nullifiers \emph{evolve} per epoch, so consensus
checks duplicates over a bounded window and permanently prunes older state;
and untrusted \emph{oblivious synchronization services} keep wallets' proofs
current without learning their transaction details (Bowe, ``Tachyon: Scaling
Zcash with Oblivious Synchronization'', 2025-04-02; Bowe--Miers, eprint
2025/2031, 2025-11-03). The intended proof layer is Ragu, a recursive-proof
(PCD) framework over the Pasta cycle with no trusted setup
(\texttt{tachyon-zcash/ragu} @ \texttt{830bbcda}). None of this is deployed;
none of it is specified at consensus rigor.

\subsection{Why a design-stage volume}\label{sec:tachyon-why}

The companion volumes document a deployed protocol: every claim in the
\emph{Ironwood Guide} cites shipped code and a final specification. Tachyon has
neither --- no protocol specification, no merged ZIP, no code on any network.
It nevertheless warrants a volume now, for two reasons. First, its standing:
of all proposals in the February 2026 NU7 sentiment poll, Tachyon drew the
most unambiguous support (\S\ref{sec:tachyon-standing}), so it is the
direction Zcash's shielded layer is most likely to take, and the design
deserves scrutiny \emph{before} it hardens into consensus rules. Second, its
volatility: the design is single-committer, actively churning, and
self-admittedly stale in places
(\S\ref{sec:tachyon-sources}); fixing its current state against commit-pinned
sources creates the baseline any later normative work can be diffed against.
The cost of writing now is stated plainly: this volume can assert nothing
normatively, and its strongest chapter is the one on open problems.

\subsection{The three-bin method}\label{sec:tachyon-bins}

Design-stage material invites a failure mode the deployed-protocol volumes
never face: prose that reads as specification but rests on a blog post. We
guard against it by classifying every claim in this volume into exactly one of
three bins, and keeping the bin visible in the prose.

\begin{definition}[Three-bin classification]\label{def:three-bins}
A claim about Tachyon is exactly one of:
\begin{enumerate}[label=(\roman*)]
\item \emph{specified} --- a normative artifact pins it down (a merged ZIP, the
  protocol specification, or deployed code); we cite the artifact, with file
  and commit;
\item \emph{designed but unspecified} --- an informal source (design book,
  eprint note, blog post, forum post, issue) describes a concrete mechanism,
  but no specification fixes it; we cite the source and its date (plus commit
  for repositories) and state what a specification would still need to pin
  down;
\item \emph{an open problem} --- no source, formal or informal, resolves it;
  where the designers themselves acknowledge the gap, we quote them.
\end{enumerate}
\end{definition}

The bins are this volume's organizing device. Each subsequent section groups
its material by bin --- \S\ref{sec:tachyon-specified} carries the specified
bin exhaustively, \S\ref{sec:tachyon-dbu} and \S\ref{sec:tachyon-open} the
lower two --- or closes each construction with a \emph{Classification}
paragraph naming its bin; a claim appearing outside such a block names its
bin inline. Claims about deployed or extant code cite
the crate and file; claims about design carry a date, because the design has
already pivoted once (\S\ref{sec:tachyon-sources}) and undated statements from
the wrong era are actively misleading.

\subsection{The source landscape}\label{sec:tachyon-sources}

Table~\ref{tab:tachyon-sources} lists the primary sources, each pinned to a
commit or date. All repository facts below were verified firsthand at the
stated commits; web sources were fetched 2026-07-08.

\begin{table}[htbp]
\centering
\small
\begin{tabular}{@{}lll@{}}
\toprule
Artifact & Pin / date & Rigor \\
\midrule
eprint 2025/2031 (Bowe--Miers)
  & rev.\ 2025-11-03            & 8-page note; no theorems \\
\texttt{tachyon-zcash/tachyon} (book, crate)
  & \texttt{26fdc165}, 2026-07-07 & design book; crypto stubbed \\
\texttt{tachyon-zcash/ragu}
  & \texttt{830bbcda}, 2026-07-05 & pre-release code; unaudited \\
\texttt{tachyon-zcash/zips} PRs \#1--\#4
  & opened 2026-06-02; open     & drafts, not upstreamed \\
\texttt{zcash/zips} (upstream)
  & \texttt{c6b70358}, 2026-07-05 & zero Tachyon ZIPs \\
Blog, forum t/50789, ZK Podcast 388
  & 2025-04-02 to 2026-01-31    & designer prose \\
\bottomrule
\end{tabular}
\caption{Primary Tachyon sources, pinned. Verification date 2026-07-08.}
\label{tab:tachyon-sources}
\end{table}

\paragraph{The eprint note.} The canonical publication is Bowe and Miers, ``A
Note on Notes: Towards Scalable Anonymous Payments via Evolving Nullifiers and
Oblivious Synchronization'' (eprint 2025/2031; received 2025-11-02, revised
2025-11-03). It is a self-described ``short note'' of eight pages: it fixes
the two central ideas and their motivation, states the fatal risks in the
authors' own words, and contains zero theorems and no formal security
definitions. It anchors bin~(ii) for the core concepts and, unusually, much of
bin~(iii): its frankest passages are the acknowledged unsolved problems.

\paragraph{The project book and crate.} The repository
\texttt{tachyon-zcash/tachyon} @ \texttt{26fdc165} (2026-07-07) carries the
most concrete design material: an mdBook elaborating nullifier derivation, the
proof tree, the anchor chain, aggregation, and a ZIP dependency map, beside
the \texttt{zcash\_tachyon} crate. Three caveats govern every citation of it.
First, it is a single-committer artifact (\texttt{git shortlog} lists only Tal
Derei) with no independent publication of its post-eprint machinery. Second,
the crate does not implement the load-bearing cryptography: its README lists
note commitments, nullifier derivation, and proof creation, merging, and
verification as ``Stubbed (pending Ragu PCD and Poseidon integration)''
(\texttt{README.md} line~49 @ \texttt{26fdc165}). Third, the book is
self-admittedly stale: issue \#121 (2026-05-26) records a pivot in the payment
design --- ``Since we're using PIR to make repeated payments flows without
OOB, there are stale parts of the book that will need to change. We're
retaining in-band secret distribution for recover from seed flows'' --- so
every statement from the out-of-band era (the 2025 blog posts, the book's key
chapter) must be dated pre-2026-05-26, and we date them throughout.

\paragraph{The ZIP landscape.} The project's ZIP program exists at two rigors,
and the split is the clearest single measure of Tachyon's maturity. The five
\emph{core} ZIPs --- Shielded Protocol, Bundle, Accumulator, Aggregator,
Oblivious Synchronization (tracking issues \#103--\#107) --- are empty stubs:
each book file carries only a tracking link and blank section headings
(\texttt{book/src/zips/} @ \texttt{26fdc165}, 216--245 bytes per stub). Four
\emph{peripheral} amendments --- to ZIP 317 (fees), ZIP 209 (pool balance),
ZIP 221 (history tree), and ZIP 248 (bundle registration) --- exist as
ZIP-0-conformant drafts, but only as PRs \#1--\#4 on the project's own
\texttt{zcash/zips} fork, open since 2026-06-02 and not upstreamed; the ZIP
248 they amend is itself unmerged upstream (open PR zcash/zips\#1156, idle
since 2026-05-26). Upstream \texttt{zcash/zips} @ \texttt{c6b70358} contains
zero Tachyon ZIPs: a repository-wide search finds only the ZIP editor list,
one author affiliation in \texttt{protocol/protocol.tex}, and one footnote
citing the founding blog post. The NU7 deployment draft
(\texttt{draft-arya-deploy-nu7.md}, Status Draft, branch \texttt{0x77190AD8},
activation heights TBD) schedules no Tachyon ZIP --- nor any other.

\paragraph{Ragu.} The proof layer lives in \texttt{tachyon-zcash/ragu}
@ \texttt{830bbcda} (2026-07-05): a working PCD pipeline over the Pasta cycle,
at workspace version 0.0.0, with no release, no crates.io publication, no
audit, and in-code security placeholders (the Fiat--Shamir domain tag is
literally \texttt{b"FIXME"}; \texttt{crates/ragu\_pcd/src/lib.rs}
lines~50--51).
We treat Ragu's code state in its own section; here it matters only as a
pinned source.

\paragraph{Informal sources.} The remaining corpus is designer prose: Bowe's
blog (founding post 2025-04-02; ``Tachyaction at a Distance'' 2025-05-15;
post-quantum post 2025-10-16, after which the index ends and the promised
deep-dive posts never appeared), the forum megathread t/50789 (28 posts,
2025-04-02 to 2026-01-31), the project site (dateless roadmap), and the ZK
Podcast 388 transcript (2026-01-21). These sources carry bin~(ii) intent and
bin~(iii) admissions, never specification. We quote them with dates. The
``two orders of magnitude'' headline is the project's own (tachyon.z.cash)
with no published derivation; the circulating TPS estimates are analyst
characterization, not project claims; no primary document derives either
(\S\ref{sec:tachyon-figures}).

\subsection{Standing}\label{sec:tachyon-standing}

Two institutional facts frame the volume. First, governance: in the Zcash
Foundation's February 2026 NU7 sentiment poll, Tachyon drew ``universal or
near-universal support across every group'' --- ``the most unambiguous signal
from the entire poll'' --- with the stated next step to ``assess technical
readiness and scope for inclusion in NU7'' (zfnd.org, 2026-02-23). Support is
not scheduling: as noted above, the NU7 deployment draft lists no ZIPs at all.
Second, editorship: Sean Bowe and Tal Derei serve as ZIP editors ``associated
with Project Tachyon'' (\texttt{zip-0000.rst} line~196 @ \texttt{c6b70358}),
and the protocol specification's author block carries the same affiliation
(\texttt{protocol/protocol.tex} line~2953). The people who would shepherd
Tachyon ZIPs through the process are the protocol's designers --- an
institutional fact we record without inference.

\subsection{What this volume can and cannot deliver}\label{sec:tachyon-expect}

In bin~(i) we can place only the standing environment: the poll outcome and
editorships above,
the adjacent network-upgrade context (NU6.2 final, NU6.3 in draft), the
out-of-band substrate ZIPs the design leans on, and Ragu's code state at
\texttt{830bbcda}. Bin~(ii) holds the entire core protocol --- evolving
nullifiers, dual-nullifier spends, the anchor chain, the proof tree,
delegation, aggregation, the pool decision --- every part of it cited to the
book at \texttt{26fdc165} as non-normative, dated, and in places already
superseded. No byte format, consensus rule, security theorem, or performance
figure appears in this volume as a project claim, because none exists at that
rigor.

The strongest chapter is therefore the one on open problems, and this is not a
concession but the point: the best-documented statements in the corpus are the
designers' own acknowledgments of what remains unsolved --- timing correlation
by synchronization services (``fatal to the entire approach''; eprint
2025/2031 \S1.3), data availability after validator pruning (``unspendable,
potentially permanently''; eprint 2025/2031 \S1.4), and recoverability, whose
entire current written resolution is one sentence in issue \#121. The events
that would obsolete parts of this volume are catalogued as revisit triggers
in \S\ref{sec:tachyon-revisit}; we revisit at any of them.

\section{The proof-carrying wallet}\label{sec:tachyon-architecture}

An Orchard wallet reconstructs its state from the public chain. To detect its
notes it trial-decrypts every published Action from its birthday height
onward, and to keep its authentication paths current it ingests every
published commitment --- its own and everyone else's --- into a pruned local
tree (the \emph{Ironwood Guide}, \S``Transmission and detection of a note:
encryption, trial decryption, and synchronisation'' and \S``Proof of
ownership of \emph{some} valid note: the commitment tree''; the deployed
mechanism is librustzcash's \textsf{shardtree}). Both costs scale with global
chain traffic, not with the wallet's own activity. Tachyon inverts the
relationship: the wallet holds a \emph{proof} of its own spendability ---
``this note entered the pool, and no nullifier of it has been published
since'' --- and advances that proof by folding in each block's deltas,
retaining nothing else about the chain. The concept is Bowe's founding post
(``Tachyon: Scaling Zcash with Oblivious Synchronization'', 2025-04-02); the
publication of record is the Bowe--Miers note (eprint 2025/2031, 2025-11-02),
eight pages with no theorems; the concrete elaboration is the project book
(tachyon-zcash/tachyon, commit \texttt{26fdc165}, 2026-07-07).

Nothing in this section is consensus-specified. Upstream zcash/zips carries
no Tachyon ZIP (repository state at commit \texttt{c6b70358}, 2026-07-05),
and the core ``Tachyon Shielded Protocol'' ZIP is an empty stub
(tachyon-zcash/tachyon issue \#103). We therefore tag every claim with its
bin: \emph{specified} (a citable artifact at spec or code rigor),
\emph{designed but unspecified} (an informal source, together with what a
specification would still need to pin down), or \emph{open problem}.

\subsection{Wallet state as an accumulation instance}\label{sec:tachyon-fold}

The state object is the \emph{spendable}, the triple
\[
(\,\cmsym,\ \nfsym_e,\ \mathsf{anchor}\,),
\]
carried as a proof header: the note's commitment, its nullifier for the
current epoch $e$, and a running commitment to the pool state up to which
absence has been proven (tachyon book, \texttt{src/proof-tree.md} and
\texttt{src/nullifiers.md} at \texttt{26fdc165}). A Tachyon note has a
distinct nullifier per epoch: from the note's trapdoor $\psi$ and the
wallet's nullifier key $\nksym$ a per-note master key
$\mathsf{mk} := \mathsf{Poseidon}_{\texttt{Tachyon-NfPrefix}}(\psi, \nksym)$
roots a GGM tree whose leaf at epoch $e$ yields $\nfsym_e$, and the
commitment
$\cmsym := \mathsf{Poseidon}_{\texttt{Tachyon-CmDerive}}(\rcm, \mathsf{pk},
v, \psi)$ digests $\psi$ and $\mathsf{pk}$ (which pins $\nksym$), so a note's
entire nullifier sequence is frozen when its commitment enters the pool.
Construction~\ref{def:tachyon-design-climb} states the derivation in full
--- the climb and leaf domains, and the working depth and arity the crate
pins (\S\ref{sec:tachyon-design-ggm}) while the book leaves them symbolic.
The proof behind the header certifies a single sentence: the note whose
commitment is $\cmsym$ was minted --- its creation stamp is folded into the
pool state --- and none of its nullifiers, from the creation epoch through
$e$, appears in the pool up to $\mathsf{anchor}$.

The fold advances this state without the wallet touching the chain. A
synchronisation service --- holding only an opaque window of the note's
future nullifiers, never the note, $\cmsym$, $\psi$, or $\mathsf{mk}$ ---
assembles an \emph{Unspent} segment: \texttt{UnspentSeed} proves one
wallet-supplied nullifier absent from one stamp's tachygram set,
\texttt{EmptyBlockUnspentSeed} covers empty blocks, \texttt{UnspentFuse}
composes contiguous segments, and \texttt{UnspentEpochFuse} crosses an epoch
boundary, splicing the completed epoch's nullifier into the segment's
elapsed history. The wallet then binds and folds: \texttt{VerifyUnspent}
proves the segment's crossed nullifiers and its tip are genuine GGM leaves of
the note identified by $\cmsym$, and \texttt{SpendableLift} consumes the
prior spendable together with the verified segment, checks commitment
equality, nullifier continuity, and anchor adjacency, and emits the advanced
spendable (book, \texttt{src/proof-tree.md} at \texttt{26fdc165}). Each step
consumes at most two prior proofs and emits one --- proof-carrying data over
the accumulation-style recursion of the Halo lineage (the \emph{Halo~2
Guide}, \S``Recursive proof composition and accumulation schemes''; the
Halo-specific device in \S``Accumulation and the Halo trick'').

The intended proof system is Ragu, a Rust PCD framework implementing a
modified Halo recursive SNARK with no trusted setup over the Pasta cycle. A
working pipeline exists in code --- application registration, \texttt{seed}
(leaf proofs), \texttt{fuse} (binary merge), \texttt{rerandomize},
\texttt{verify}, with passing multi-node PCD-tree integration tests
(\texttt{crates/ragu\_pcd/src/\{lib.rs,fuse/mod.rs,verify.rs\}} at commit
\texttt{830bbcda}, 2026-07-05). The same code is explicit that it is
pre-release: the \texttt{Proof} struct carries full polynomials and
verification is linear-time (no inner-product opening argument or decider
exists anywhere in the workspace), the Fiat--Shamir domain tag is the
literal \texttt{b"FIXME"} (\texttt{crates/ragu\_pcd/src/lib.rs:50--51}), and
registry binding carries an in-code \texttt{FIXME(security)}. Downstream,
the \texttt{zcash\_tachyon} crate registers all nineteen proof-tree steps but
compiles only against Ragu's \texttt{mock} feature, with note commitments,
nullifier derivation, and proof creation/merge/verification ``Stubbed
(pending Ragu PCD and Poseidon integration)'' (tachyon \texttt{README.md} at
\texttt{26fdc165}).

\paragraph{Classification.} The Ragu and \texttt{zcash\_tachyon} code facts
are \emph{specified} in the only sense available before a spec: verifiable
source at a pinned commit. The wallet-state design --- the headers, the
nineteen steps, the fold --- is \emph{designed but unspecified}: it exists in
the project book and nowhere at ZIP rigor, and a specification would need to
pin the header encodings, the per-step relations, the Poseidon instances and
domain tags, the GGM depth and arity, and the epoch length. The book's
soundness section is prose --- no formal definitions, theorems, or
machine-checked statements, and the eprint contains none either --- so
recursive soundness of the tree, and whether Ragu's linear-time verifier can
be carried to a succinct decider at the sizes Tachyon needs, are \emph{open
problems} (ragu \texttt{crates/ragu\_pcd/src/verify.rs:122} records three
missing verifier checks; the performance figures in the Ragu book are
self-labelled rough estimates).

\subsection{What the fold removes}\label{sec:tachyon-removes}

Two cost centres of the deployed protocol disappear; both removals are
properties of the \emph{designed-but-unspecified} core above and inherit its
bin except where noted.

\paragraph{Trial decryption of the chain.} Nothing on today's chain marks
which Action pays a wallet, so detection is exhaustive: the wallet
trial-decrypts every Orchard output from its birthday height, and
restore-from-seed replays that scan over the entire chain since the birthday
(\emph{Ironwood Guide}, \S``Transmission and detection of a note: encryption,
trial decryption, and synchronisation''). Under Tachyon the wallet never
scans. Detection and delivery move to the payment channel
(\S\ref{sec:tachyon-payments}), and the absence-tracking that
synchronisation actually requires is delegated to a service that works only
on the opaque nullifier windows of \S\ref{sec:tachyon-fold}: the eprint
asserts services are ``fully oblivious to their clients' transaction
details'' (eprint 2025/2031), and the book's role table assigns the service
the \texttt{Unspent} steps beside the note-free anchor-chain segments, and
none that touch the note (\texttt{src/proof-tree.md} at \texttt{26fdc165}).
Groundwork exists at ZIP
rigor: ZIP~2004, ``Remove the dependency of consensus on note encryption''
(Status: Draft at \texttt{c6b70358}), makes the consensus layer independent
of the ciphertexts a delivery redesign would drop, though its
\texttt{ephemeralKey} handling is an in-file TODO. No formal definition of
``oblivious'' exists --- no threat model and no security definition --- which
we treat as an \emph{open problem} (the trust model,
\S\ref{sec:tachyon-dbu-oss}; the timing-correlation attack the eprint calls
potentially fatal, \S\ref{sec:tachyon-open-timing}).
The aggregate cost claim is likewise unproven: the designer's estimate that
total service effort is ``asymptotically similar'' to today's validators is
a forum statement (t/50789 \#28, 2025-11-21) with no published complexity
analysis, and no project benchmark exists --- circulating throughput and
size figures are analyst estimates (Messari; CoinDesk Research, 2026-06-30),
not project figures.

\paragraph{Per-note witness maintenance.} Today the wallet appends every
published $\cmx$ to a pruned local shard tree, marks the positions of its own
notes, stores a checkpoint per block, and rebuilds each authentication path
on demand from retained subtree roots (\emph{Ironwood Guide}, \S``Proof of
ownership of \emph{some} valid note: the commitment tree''; the design of
librustzcash's \textsf{shardtree}). Tachyon has no note commitment tree.
Pool state is a hash chain (\S\ref{sec:tachyon-validators}); membership of
the note's creation is proven once, when \texttt{SpendableInit} checks
$\cmsym$ against the creation stamp's tachygram set; and the lineage
thereafter carries no path --- there are no siblings to refresh as others
append, because there is no tree (book, \texttt{src/proof-tree.md},
\texttt{src/anchor.md} at \texttt{26fdc165}). What remains is the forward
obligation: the wallet must keep proving absence up to a recent anchor
before it can spend, work proportional to its own note count and elapsed
epochs rather than to chain traffic, and delegable as above. \emph{Designed
but unspecified} throughout: the creation-stamp membership argument and the
anchor-rooting rule exist only in the book.

\subsection{Validator-side consequences}\label{sec:tachyon-validators}

A node today maintains two unbounded structures: the append-only commitment
tree, from which it recognises anchors, and the ever-growing nullifier set,
against which it checks each Action's $\nfsym_{\mathrm{old}}$ --- the one
check that cannot sit inside a proof, because it concerns global history
(\emph{Ironwood Guide}, \S``What a node verifies, without ever observing the
note''). The Tachyon design replaces both with one running Poseidon hash
chain and bounds the duplicate check to a recent window.

Each stamp accepted into a block advances the pool state by
\[
\mathsf{anchor}' :=
  \mathsf{Poseidon}_{\texttt{Tachyon-StampFld}}(\mathsf{anchor},\ x,\ y),
\]
where $(x, y)$ are the coordinates of the stamp's tachygram-set commitment;
an empty block advances it by
$\mathsf{Poseidon}_{\texttt{Tachyon-EmptyBlk}}(\mathsf{anchor})$; and the
crossing into epoch $e + 1$ by
$\mathsf{Poseidon}_{\texttt{Tachyon-EpochStp}}(\mathsf{anchor},\ e + 1)$,
the only domain that absorbs the epoch index (book, \texttt{src/anchor.md}
at \texttt{26fdc165}). Anchors are the end-of-block states of this chain; a
spendable or stamp proves continuity from one anchor to a later one rather
than membership in a tree.

The double-spend check becomes a windowed duplicate check: consensus checks
every tachygram of a published stamp for duplicates over every block of the
current and the immediately preceding epoch, rejects on a hit, and may
permanently prune tachygrams older than that window (book,
\texttt{src/tachygrams.md} at \texttt{26fdc165}). Two epochs suffice in the
design because a stamp's anchor may have been lifted within its epoch (so
the whole anchor epoch is scanned, not merely blocks after the anchor), and
because a spend publishes the note's nullifiers for both the anchor epoch
$e$ and epoch $e + 1$ --- the next-epoch nullifier published in $e$ is
exactly what a second spend of the same note would publish in $e + 1$
(book, \texttt{src/nullifiers.md} at \texttt{26fdc165}).

\paragraph{Classification.} The window rule and prunability are
\emph{designed but unspecified}: they appear in the book only; the eprint
states just that ``validators must still check for duplicates''; no
consensus ZIP drafts the rule. A specification would need the epoch length
--- the parameter the entire pruning model hinges on, which the project
itself marks as requiring empirical analysis (tachyon-zcash/tachyon issue
\#79) --- and the normative force of end-of-block anchoring, on which the
book's own source wavers between ``may'', ``should'', and ``must'' in an
inline comment (\texttt{src/anchor.md}). The consequence of
pruning is an \emph{open problem}: what validators discard, wallets and
services still need, and the eprint concedes that the sudden shutdown of
syncing services ``could render the funds of all their dependent users
unspendable, potentially permanently'' (eprint 2025/2031, \S1.4); no
data-availability or incentive design exists (\S\ref{sec:tachyon-open-da}).

\subsection{Payment delivery}\label{sec:tachyon-payments}

Removing detection-by-trial-decryption reopens the question the deployed
protocol answers with in-band secret distribution (\emph{Ironwood Guide},
\S``Transmission and detection of a note: encryption, trial decryption, and
synchronisation''): how does the recipient learn of the note? The project's
answer changed in May 2026. We present both stages and date every statement;
all quotations are verbatim from the cited artifacts.

\paragraph{The out-of-band era (all statements predate 2026-05-26).} The
design of ``Tachyaction at a Distance'' (Bowe, 2025-05-15) made Tachyon
``out-of-band payments for the first time in a Zcash shielded protocol'':
the sender delivers the note plaintext to the recipient over a channel
outside the chain, and the on-chain ciphertext disappears. The
simplification cascades through the key hierarchy --- key diversification,
viewing keys, and payment addresses are removed; the note drops the
diversifier and $\rho$; the diversified transmission key $\pkd$ gives way to
a payment key deterministic per spending key
(\S\ref{sec:tachyon-design-keys}, which owns the construction; the book's
\texttt{src/keys.md} opens with the in-file warning ``this is
significantly outdated''). Unlinkability becomes the wallet layer's problem,
deferred to the URI substrate: ZIP~321 payment-request URIs and ZIP~324
URI-encapsulated payments. That substrate is \emph{specified} only in part.
ZIP~321 is Status: Active. ZIP~324 is Status: Draft, predates Tachyon and
mentions it nowhere, declares that ``It is outside the scope of this
proposal to establish a secure communication channel'', derives its
\textsf{rseed} under a domain separator that is the literal
\texttt{0xFIXME}, is Sapling-only, and ends with a ``Publication Blockers''
section (zips \texttt{zip-0324.rst} at \texttt{c6b70358}). ZIP~231 memo
bundles (Status: Draft) decouple memo data from individual outputs at full
cryptographic rigor, but their v6 carrier was orphaned when ZIP~230 was
withdrawn (zips \texttt{zip-0231.md}, \texttt{zip-0230.rst} at
\texttt{c6b70358}). The out-of-band payment protocol itself was never
specified at any rigor; the project roadmap listed two unreconciled delivery
approaches (tachyon.z.cash/roadmap/, fetched 2026-07-08). The founding post
had already conceded the cost: without in-band distribution ``the user
cannot rely on the blockchain as a storage mechanism for recovering their
funds from a seed phrase'' (Bowe, 2025-04-02).

\paragraph{The pivot (2026-05-26).} Issue \#121 of tachyon-zcash/tachyon,
opened 2026-05-26 and still open as of 2026-07-08, replaces the out-of-band
story in two sentences, which we quote whole:
``Since we're using PIR to make repeated payments flows without OOB, there
are stale parts of the book that will need to change. We're retaining
in-band secret distribution for recover from seed flows.'' Repeated-payment
flows are thus to use private information retrieval rather than
sender-to-recipient out-of-band delivery --- the roadmap adds a post-quantum
key exchange for the channel --- while in-band secret distribution returns,
scoped to recover-from-seed and restoring exactly the property the founding
post conceded was lost. The protocol is ``developed independently in
collaboration with'' the team and has no published specification; its
public trace is a roadmap paragraph and podcast/whiteboard discussion
(tachyon.z.cash/roadmap/; ZK Podcast 388, 2026-01-21). \emph{Designed but
unspecified} at its thinnest: a specification would need the PIR scheme and
its database and server model, the key-exchange construction, the
recover-from-seed format (which secrets stay in-band, under which keys, at
what fee treatment), and the consequences for the key-hierarchy removals
above --- the book is flagged stale against the pivot by issue \#121
itself, so which of \texttt{src/keys.md}'s removals survive is
undetermined. The post-quantum claim also needs scoping: Bowe's ``Zcash and
Quantum Computers'' (2025-10-16) limits Tachyon's quantum benefit to privacy
--- the harvest-now-decrypt-later threat --- while ECC-based proof soundness
remains quantum-vulnerable; whether retained in-band recovery ciphertexts
reuse the ECC key agreement, and so retain its harvest-now exposure, is
undetermined. Recoverability at large remains an \emph{open problem} by the
designer's own account: ``Getting this right will be challenging, but we
simply have no choice but to do it'' (forum t/50789 \#22, 2025-11-04); on ZK
Podcast 388
(2026-01-21) the recovery infrastructure was still ``a can of worms''. The
single sentence of issue \#121 is the entire written resolution to date.

\section{The design as published}\label{sec:tachyon-design}

The preceding section reconstructed the architecture; this one fixes the
design's object model as the project book publishes it: the mdBook of
\texttt{tachyon-zcash/tachyon} @ \texttt{26fdc165} (2026-07-07,
single-committer; \S\ref{sec:tachyon-sources}), the design's most complete
public statement. Each subsection states what the book fixes --- exact
formulas, fields, and domain strings --- names the Orchard analogue in one
contrast, and closes with the object's bin under
Definition~\ref{def:three-bins} and what a specification must still pin
down; mechanisms the architecture section constructs and gaps the status
sections classify are cross-referenced, not restated. Every claim here is
\emph{designed but unspecified} by default. Where the
\texttt{zcash\_tachyon} crate at the same commit realizes an object, the
realization is \emph{specified} in the only sense available before a spec
--- verifiable source at a pinned commit (\S\ref{sec:tachyon-crate}) ---
with the README split as the boundary: keys, value commitments, signatures,
and bundle construction ``implemented and tested''; note commitments,
nullifier derivation, proofs, and stamp compression ``Stubbed''. One caveat
cuts across every Poseidon-derived object: the crate routes all Poseidon
calls through Ragu's mock sponge (\texttt{src/digest/poseidon.rs}), so
derivation \emph{structure} and domain tags are code-anchored while the
outputs are not real Poseidon. (The book's $\F_p$, $\F_q$ are the series'
$\F_{p_{\Pallas}}$, $\F_{p_{\Vesta}}$.)

Staleness is stated once, here. The book post-dates the May-2026 payment
pivot (issue \#121; \S\ref{sec:tachyon-sources}), yet two chapters carry
the in-file marker \texttt{<!--\ todo: this is significantly outdated -->}
(\texttt{src/keys.md} line~3, whole chapter; \texttt{src/authorization.md}
line~228, its Detailed Sequence), and the key chapter's out-of-band claims
are pre-pivot statements, dated pre-2026-05-26
(\S\ref{sec:tachyon-payments}, \S\ref{sec:tachyon-dbu-pir}); we document
the book as published and do not relitigate per object. Eight
\texttt{src/SUMMARY.md} entries carry no file: the four section parents
(Data Model, Cryptographic Primitives, Protocol, Network Roles --- the
last has a filed Overview child), the three network-role pages, and the
testnet chapter. One terminological drift: the bundle chapter's
\emph{tachystamp} is the \emph{stamp} of the proof-tree chapter and the
crate; we write \emph{stamp}.

\subsection{Bundles and tachygrams}\label{sec:tachyon-design-bundle}

\begin{definition}[Tachyon bundle objects]\label{def:tachyon-design-bundle}
A \emph{tachygram} is a 32-byte field element representing either a
nullifier or a note commitment; consensus does not distinguish the two. A
\emph{tachyaction} indistinguishably represents the creation or
destruction of a note --- a commitment tachygram proves creation, a
nullifier tachygram destruction --- and is cryptographically bound to
one tachygram that it does \emph{not} contain. It contains
$\cvsym$, a 32-byte homomorphic commitment to the created or destroyed
value; $\rk$, a 32-byte public key bound to one tachygram; and
$\mathsf{sig}$, a 64-byte RedPallas signature by $\rk$ over the
transaction sighash. A \emph{stamp} contains an anchor, the recursive
proof (which may be aggregated), and the covered actions' tachygrams; the
proof establishes that each tachygram creates a new note or destroys an
existing one, that tachygrams are correctly bound to action keys, and that
the action balance effect matches the pool balance effect. A \emph{bundle}
collects the tachyactions, the integer net pool effect
$\mathsf{value\_balance}$, a binding signature over the same sighash as
the action signatures, and the stamp (book, \texttt{src/bundle.md} @
\texttt{26fdc165}).
\end{definition}

The covering proof witnesses a commitment to the \emph{unordered set} of
the stamp's tachygrams as a root polynomial,
$\mathsf{tachygram\_acc}(X) = \prod_i (X - \mathsf{tg}_i)$, and this set is
the atomic unit that folds into pool state to contribute to the anchor
(book, \texttt{src/tachygrams.md} @ \texttt{26fdc165}); the fold and the
two-epoch duplicate scan the set feeds are constructed at
\S\ref{sec:tachyon-validators}.

\paragraph{Wire format.} The first byte, \texttt{tachyonBundleState},
selects one of three states: $\mathtt{0x00}$ non-tachyon (no bundle),
$\mathtt{0x01}$ stamped, $\mathtt{0x02}$ stripped; the body layout (value
balance as i64, the $(\cvsym, \rk)$ action vector, per-action and binding
signatures) is fixed only in the crate's doc tables
(\texttt{crates/tachyon/src/bundle/mod.rs}). The stamp trailer (state
$\mathtt{0x01}$) carries \texttt{cActionsTachyon} (32 bytes, a
``compressed Pedersen commitment'' to the stamp's merged action-digest
set; \S\ref{sec:tachyon-design-agg}), \texttt{anchorTachyon} (32 bytes),
\texttt{nTachygrams} (compactsize), \texttt{vTachygrams} ($32n$ bytes),
and \texttt{proofTachyon}, a \texttt{PROOF\_SIZE} blob --- ``serialized
proof of fixed size''. The stripped trailer (state $\mathtt{0x02}$)
carries only \texttt{tachyonAggregateId}, the 64-byte nonzero
\texttt{wtxid} of the covering aggregate: the all-zero \texttt{wtxid},
naming no aggregate, is rejected on read even for a stripped innocent, and
the stripped trailer carries no \texttt{cActionsTachyon} --- that field
rides on the stamp and strips away when a bundle becomes an adjunct.

The Orchard contrast: an Orchard Action publishes typed, distinguishable
fields ($\cvsym^{\mathrm{net}}$, $\nfsym^{\mathrm{old}}$, $\rk$, $\cmx$)
plus note ciphertexts, with one per-transaction Halo~2 proof (the
\emph{Ironwood Guide}, \S``What a node verifies, without ever observing the
note''); a tachyaction is single-effect, ciphertext-free, does not carry
its tachygram, and its proof is strippable and aggregatable across
transactions --- a lifecycle with no Orchard analogue.

\paragraph{Classification.} \emph{Designed but unspecified}, with the wire
format's structure specified-as-code: the crate's doc tables mirror the
book's byte-for-byte and the nonzero-\texttt{wtxid} rule is enforced on
read (\texttt{bundle/mod.rs}, \texttt{stamp/mod.rs}). A specification must
still pin \texttt{PROOF\_SIZE}, which no number in the book fixes (the
crate defers to Ragu's \texttt{PROOF\_SIZE\_COMPRESSED}), and the
point-compression convention, fixed only informally per context (the
trailer compresses; the anchor absorbs complete coordinates,
\S\ref{sec:tachyon-design-tree}).

\subsection{The key hierarchy}\label{sec:tachyon-design-keys}

\begin{definition}[Tachyon key hierarchy]\label{def:tachyon-design-keys}
The spending key $\mathsf{sk}$ is raw 32-byte entropy, matching Orchard's
representation and preserving the full 256-bit key space. With
$\mathsf{PRFexpand}_{\mathsf{sk}}(t) =
\text{BLAKE2b-512}(\texttt{Zcash\_ExpandSeed},\ \mathsf{sk} \,\|\, t)$,
the book derives
\[
\ask = \mathrm{ToScalar}\bigl(\mathsf{PRFexpand}_{\mathsf{sk}}
  ([\mathtt{0x21}])\bigr), \qquad
\nksym = \mathrm{ToBase}\bigl(\mathsf{PRFexpand}_{\mathsf{sk}}
  ([\mathtt{0x22}])\bigr),
\]
a long-lived scalar in $\F_{p_{\Vesta}}$ and a base-field element in
$\F_{p_{\Pallas}}$, with $\aksym = [\ask]\,\mathcal{G}$; the payment key is
\[
\mathsf{pk} = \mathsf{Poseidon}(\texttt{PK\_DOMAIN},\ \aksym_x,\ \nksym),
\]
with $\aksym_x$ the $x$-coordinate of $\aksym$ --- the chapter leaves
\texttt{PK\_DOMAIN} symbolic; only the domain-separator table fixes it as
\texttt{Tachyon-PkDerive} (\S\ref{sec:tachyon-design-txid}). The
\emph{proof authorizing key} $\mathsf{pak} = (\aksym, \nksym)$ enables
proof construction without spend authority; the per-note delegation bundle
$(\aksym, \mathsf{mk})$ restricts a prover to one note, all epochs, where
$\mathsf{mk}$ is the note's GGM master key
(\S\ref{sec:tachyon-design-ggm}). (Book, \texttt{src/keys.md} @
\texttt{26fdc165}.)
\end{definition}

\paragraph{Sign normalization and tiers.} RedPallas requires
$\aksym = [\ask]\,\mathcal{G}$ to have $\tilde{y} = 0$ (bit 255 of the
compressed encoding); if $\tilde{y}(\aksym) = 1$ the derivation negates
$\ask$ ($[-\ask]\,\mathcal{G} = -[\ask]\,\mathcal{G}$), once, at derivation
time (\texttt{src/keys.md}; \texttt{crates/tachyon/src/keys/private.rs}
lines 83--98). The key $\ask$ cannot sign directly and $\aksym$ cannot
verify directly: both are randomized per action into $\rsk$ and $\rk$
(\S\ref{sec:tachyon-design-auth}), so each on-chain $\rk$ is unlinkable to
$\aksym$; the book names the tiers private scalar $\to$ circuit witness
$\to$ public on-chain (\texttt{SpendAuthorizingKey},
\texttt{SpendValidatingKey}, \texttt{ActionVerificationKey}). The payment
key is \emph{deterministic per spending key} --- every note from one
$\mathsf{sk}$ shares one $\mathsf{pk}$, no per-note diversification --- and
since $\mathsf{pk}$ is committed in $\cmsym$, accumulator membership
transitively pins $\aksym$ and $\nksym$: a wrong $\nksym$ gives a wrong
$\mathsf{pk}$, a wrong $\cmsym$, and inclusion fails.

The Orchard contrast: the root $\mathsf{sk}$, the $\mathsf{PRFexpand}$
construction, the $\ask \to \aksym$ derivation with the same sign
normalization, and RedPallas re-randomization survive (the \emph{Orchard
Guide}, \S``Keys: capabilities for spending and for viewing''; tag bytes
differ, $\mathtt{0x06}$/$\mathtt{0x07}$/$\mathtt{0x08}$ against
$\mathtt{0x21}$/$\mathtt{0x22}$), while $\rivk$, $\fvk = (\aksym, \nksym,
\rivk)$, $\ivk$, $\dk$, $\ovk$, the FF1-AES256 diversifier $d$, $\gd$, and
$\pkd$ are removed, and $\mathsf{pk}$, $\mathsf{pak}$, and the per-note
$(\aksym, \mathsf{mk})$ bundle are new with no Orchard counterpart.

\paragraph{Classification.} The derivations and normalization are
\emph{specified} as tested code in the crate's key module, with
$\mathsf{pk}$ computed through the mock sponge
(\texttt{keys/private.rs}). The intent around $\mathsf{pk}$ ---
unlinkability deferred out-of-band --- is the pre-pivot claim re-examined
at \S\ref{sec:tachyon-dbu-pir}. A specification must pin
\texttt{PK\_DOMAIN} normatively, the encoding of $\aksym_x$, and --- after
the pivot --- which key-hierarchy removals survive.

\subsection{Notes and commitments}\label{sec:tachyon-design-notes}

\begin{definition}[Tachyon note]\label{def:tachyon-design-note}
A note carries exactly four fields: $\mathsf{pk} \in \F_{p_{\Pallas}}$, the
recipient's payment key; $v$, a u64 value in zatoshi with
$1 \leq v \leq 2.1 \times 10^{15}$; $\psi \in \F_{p_{\Pallas}}$, the
nullifier trapdoor; and $\rcm \in \F_{p_{\Pallas}}$, the commitment
trapdoor. Zero-value notes are forbidden. The commitment binds all four,
\[
\cmsym = \mathsf{Poseidon}_{\texttt{Tachyon-CmDerive}}
  (\rcm,\ \mathsf{pk},\ v,\ \psi),
\]
and is the published tachygram of an output action; a spend action
publishes \emph{two} nullifiers, at the present epoch $e$ and the next
epoch $e+1$, as its tachygrams (book, \texttt{src/notes.md} @
\texttt{26fdc165}; the commitment already appears in
\S\ref{sec:tachyon-fold}, the dual-nullifier rationale in
\S\ref{sec:tachyon-validators}).
\end{definition}

The Orchard contrast: an Orchard note is a 5-tuple
$(\mathsf{addr} = (d, \pkd), v, \rho, \psi, \rcm)$ under a perfectly
hiding Sinsemilla $\NoteCommit$, with $v \in [0, 2^{64})$ (the
\emph{Ironwood Guide}, \S``Representation of a note: notes and note
commitments'') and load-bearing zero-value dummies (\emph{Ironwood Guide},
\S``The Action statement, formally''); Tachyon drops the address and
$\rho$ entirely --- ``no diversifier, no rho, no unique value for faerie
gold defense'' since ``the nullifier construction doesn't require global
uniqueness'' (\texttt{crates/tachyon/src/note.rs}) --- so where Orchard
chains $\rho^{\mathrm{new}} := \nfsym^{\mathrm{old}}$ (\emph{Orchard
Guide}, \S``Prevention of double-spending: nullifiers''), $\psi$ reuse
here harms only the owner (\S\ref{sec:tachyon-dbu-nf}).

\paragraph{Classification.} \emph{Designed but unspecified}, and the
commitment scheme is the book's clearest internal churn marker: the note
and domain-separator chapters fix Poseidon under \texttt{Tachyon-CmDerive},
while the crate's module doc calls the scheme ``TBD'' --- ``(e.g.\
Sinsemilla, Poseidon)'' pending what ``is efficient inside Ragu circuits''
(\texttt{note.rs}; cited at \S\ref{sec:tachyon-dbu-nf}). A specification
must pin the scheme and the field encodings; the $\psi$-uniqueness
responsibility gap is classified at \S\ref{sec:tachyon-dbu-nf}.

\subsection{Epochs and the GGM nullifier tree}\label{sec:tachyon-design-ggm}

The derivation's endpoints --- the master key
$\mathsf{mk} = \mathsf{Poseidon}_{\texttt{Tachyon-NfPrefix}}(\psi, \nksym)$,
the per-epoch leaf nullifier, the delegation window $\delta$ with its
sentinel, and the spendable lineage --- are constructed at
\S\ref{sec:tachyon-fold} and \S\ref{sec:tachyon-dbu-nf}; none of that
machinery appears in the eprint note, whose evolving-nullifier sketch is a
flat PRF still taking $\rho$ as input (quoted at
\S\ref{sec:tachyon-dbu-nf}).
What the book fixes beyond those statements is the climb itself.

\begin{construction}[GGM climb]\label{def:tachyon-design-climb}
The tree has depth $D$ and arity $A$, covering leaf epochs
$e \in [0,\ A^D - 1]$; a node at depth $d$ covers a contiguous range of
$A^{D-d}$ consecutive epochs, and a node may only climb toward the leaves,
never back toward the root --- the property that makes absence-proof
delegation trustless. To reach the leaf at epoch $e$, decompose $e$ into
$d \in [0 \ldots D]$ base-$A$ direction chunks, most significant first, and
climb one step per chunk:
\[
\mathsf{KDF}^{\mathsf{climb}}_\psi(e, 0) =
  \mathsf{KDF}^{\mathsf{root}}_\psi(\nksym) = \mathsf{mk},
\]
\[
\mathsf{KDF}^{\mathsf{climb}}_\psi(e, d) =
  \mathsf{Poseidon}_{\texttt{Tachyon-NfPrefix}}\!\left(
    \mathsf{KDF}^{\mathsf{climb}}_\psi(e, d-1),\
    \left\lfloor e / A^{D-d} \right\rfloor \bmod A
  \right).
\]
The value at $\mathsf{KDF}^{\mathsf{climb}}_\psi(e, D)$ is the leaf key for
epoch $e$; the final \texttt{Tachyon-NfDerive} hash from leaf key to
nullifier is constructed at \S\ref{sec:tachyon-dbu-nf} (book,
\texttt{src/nullifiers.md} @ \texttt{26fdc165}).
\end{construction}

\paragraph{Parameters.} The book leaves $D$ and $A$ symbolic everywhere.
The crate pins working values: $\texttt{GGM\_TREE\_ARITY} = 4$,
$\texttt{GGM\_CHUNK\_SIZE} = 2$ bits, $\texttt{GGM\_TREE\_DEPTH} = 10$
under the invariant $A^D = \texttt{GGM\_MAX\_INDEX} + 1$ with
$\texttt{GGM\_MAX\_INDEX} = \texttt{EPOCH\_MAX} = 2^{20} - 1$ ($4^{10} =
2^{20}$ epochs tile the index space exactly), and $\texttt{EPOCH\_SIZE} =
2^{12} = 4096$ blocks, via a shift constant that drops to $4$ under
\texttt{cfg(test)} --- placeholder values, not decisions
(\texttt{crates/tachyon/src/keys/ggm.rs}, \texttt{src/constants.rs}).
Nothing in the book maps the epoch index to block height; the epoch length
exists only in code, and its tuning is the open issue of
\S\ref{sec:tachyon-dbu-epoch}.

The Orchard contrast: Orchard derives \emph{one} nullifier per note,
$\nfsym = \Extract\bigl([F_{\nksym}(\rho) + \psi]\,\mathcal{G}^{\nfsym} +
\cmsym\bigr)$ --- a Poseidon PRF plus a curve point --- checked against the
global, ever-growing nullifier set (the \emph{Ironwood Guide},
\S``Prevention of double-spending: nullifiers''); Tachyon replaces it with
a per-epoch sequence, pure Poseidon with no curve point and no
$\Extract$, checked only over the two-epoch window of
\S\ref{sec:tachyon-validators}.

\paragraph{Classification.} \emph{Designed but unspecified}; the tree walk
is code-anchored in structure (\texttt{keys/ggm.rs}) but its cryptography
sits on the stubbed side of the README split. A specification must pin
$D$, $A$, the epoch length and activation semantics
(\S\ref{sec:tachyon-dbu-epoch}), and the chunk encoding; the tip-nullifier
exposure of the delegation flow is the open problem of
\S\ref{sec:tachyon-open-tip}.

\subsection{Anchors, stamps, and the proof tree}\label{sec:tachyon-design-tree}

The three anchor-chain formulas are constructed at
\S\ref{sec:tachyon-validators}; the unresolved normative force of
end-of-block anchoring is classified at \S\ref{sec:tachyon-dbu-rollover}.
The book adds two statements worth fixing here: intra-block states exist
between stamp absorptions, and a proof such as \texttt{SpendableInit}
``will produce a header that is likely at an intra-block state, and should
be lifted to the end of a block by proving state continuity''
(\texttt{src/anchor.md}; that the epoch step is the only domain absorbing
the epoch index is already fixed at \S\ref{sec:tachyon-validators} --- all
cross-epoch binding rides on it); and absorbing
the \emph{complete} $(x, y)$ coordinates of the stamp's tachygram-set
commitment, rather than a compressed encoding, makes the binding
unambiguous --- the crate sharpens this to a Vesta point whose
$\F_{p_{\Vesta}}$ coordinates are each split into two 128-bit
little-endian limbs, a six-element Poseidon absorption
(\texttt{digest/poseidon.rs}; the sequence, tachygram-set, and action-set
commitments all live on Vesta, \texttt{primitives/seq.rs},
\texttt{primitives/sets.rs}).

\paragraph{The step graph.} Each proof step accepts arbitrary witness
inputs and up to two PCD inputs, and emits a new PCD. The book fixes
exactly nineteen steps over eight header types
(Table~\ref{tab:tachyon-design-steps}) and a role matrix: the wallet runs
every
step touching the note's commitment or master key; the sync service runs
the anchor and \texttt{Unspent} steps over wallet-shared nullifier values,
never note material (\S\ref{sec:tachyon-dbu-oss}); the aggregator touches
only published \texttt{StampHeader}s and the public anchor chain --- the
three anchor-chain steps to build lift segments, then \texttt{StampLift}
and \texttt{MergeStamp} (book, \texttt{src/proof-tree.md} @
\texttt{26fdc165}).

\begin{table}[htbp]
\centering
\small
\begin{tabular}{@{}lp{10.2cm}@{}}
\toprule
Header & Fields \\
\midrule
\texttt{AnchorChain} & \texttt{(start, end)} \\
\texttt{Unspent} & \texttt{(anchor\_prev, (epoch\_start, nf\_start),
  elapsed, (epoch\_end, nf\_end), anchor\_last)} \\
\texttt{VerifiedUnspent} & \texttt{(cm, anchor\_prev, (epoch\_start,
  nf\_start), (epoch\_end, nf\_end), anchor\_last)} \\
\texttt{NfPrefixHeader} & \texttt{(cm, node, depth, index)} \\
\texttt{NullifierHeader} & \texttt{(cm, (epoch\_start, nf\_start),
  nf\_seq\_commit, (epoch\_end, nf\_end))} \\
\texttt{SpendableHeader} & \texttt{(cm, present\_nf, anchor)} \\
\texttt{SpendHeader} & \texttt{(cm, (cv, rk), present\_nf, anchor)} \\
\texttt{StampHeader} & \texttt{(action\_commit, tachygram\_commit,
  anchor)} \\
\bottomrule
\end{tabular}

\medskip

\footnotesize
\setlength{\tabcolsep}{4pt}
\begin{tabular}{@{}lllp{3.9cm}l@{}}
\toprule
Step & Left & Right & Witness & Output \\
\midrule
\texttt{AnchorSeed} & --- & --- & \texttt{start, stamp\_commit}
  & \texttt{AnchorChain} \\
\texttt{EmptyBlockSeed} & --- & --- & \texttt{start}
  & \texttt{AnchorChain} \\
\texttt{AnchorFuse} & \texttt{AnchorChain} & \texttt{AnchorChain} & ---
  & \texttt{AnchorChain} \\
\texttt{UnspentSeed} & --- & --- & \texttt{anchor\_prev, (epoch, nf),
  stamp\_tg\_set} & \texttt{Unspent} \\
\texttt{EmptyBlockUnspentSeed} & --- & --- & \texttt{anchor\_prev,
  (epoch, nf)} & \texttt{Unspent} \\
\texttt{UnspentFuse} & \texttt{Unspent} & \texttt{Unspent}
  & \texttt{left\_elapsed\_seq, combined\_elapsed\_seq,
  right\_elapsed\_seq} & \texttt{Unspent} \\
\texttt{UnspentEpochFuse} & \texttt{Unspent} & \texttt{Unspent}
  & \texttt{left\_elapsed\_seq, combined\_elapsed\_seq,
  right\_elapsed\_seq} & \texttt{Unspent} \\
\texttt{VerifyUnspent} & \texttt{Unspent} & \texttt{NullifierHeader}
  & \texttt{elapsed\_seq, nf\_seq} & \texttt{VerifiedUnspent} \\
\texttt{NfMasterSeed} & --- & --- & \texttt{note, pak}
  & \texttt{NfPrefixHeader} \\
\texttt{NfPrefixStep} & \texttt{NfPrefixHeader} & --- & \texttt{chunk}
  & \texttt{NfPrefixHeader} \\
\texttt{NullifierStep} & \texttt{NfPrefixHeader} & --- & ---
  & \texttt{NullifierHeader} \\
\texttt{NullifierFuse} & \texttt{NullifierHeader} & \texttt{NullifierHeader}
  & \texttt{left\_seq, merged\_seq, right\_seq}
  & \texttt{NullifierHeader} \\
\texttt{SpendableInit} & \texttt{AnchorChain} & \texttt{NullifierHeader}
  & \texttt{(pre\_epoch\_anchor, pre\_cm\_anchor), creation\_set,
  present\_nf} & \texttt{SpendableHeader} \\
\texttt{SpendableLift} & \texttt{SpendableHeader} & \texttt{VerifiedUnspent}
  & --- & \texttt{SpendableHeader} \\
\texttt{SpendBind} & \texttt{SpendableHeader} & --- & \texttt{note, rcv,
  alpha, pak} & \texttt{SpendHeader} \\
\texttt{OutputStamp} & --- & --- & \texttt{rcv, alpha, note, anchor}
  & \texttt{StampHeader} \\
\texttt{SpendStamp} & \texttt{SpendHeader} & \texttt{NullifierHeader}
  & \texttt{nf\_next} & \texttt{StampHeader} \\
\texttt{MergeStamp} & \texttt{StampHeader} & \texttt{StampHeader}
  & \texttt{(action\_set, tachygram\_set)} $\times$ \texttt{left, merged,
  right} & \texttt{StampHeader} \\
\texttt{StampLift} & \texttt{StampHeader} & \texttt{AnchorChain} & ---
  & \texttt{StampHeader} \\
\bottomrule
\end{tabular}
\caption{The eight proof-tree headers (top) and nineteen steps (bottom),
fields verbatim (book, \texttt{src/proof-tree.md} @ \texttt{26fdc165}).
All nineteen steps are registered in the crate against Ragu's mock
(\texttt{crates/tachyon/src/stamp/proof/mod.rs}).}
\label{tab:tachyon-design-steps}
\end{table}

\begin{construction}[Sentinel-terminated sequence composition]
\label{def:tachyon-design-splice}
An \texttt{Unspent}'s \texttt{elapsed} holds one nullifier coefficient per
epoch-boundary crossing in its span, forward-chronological, terminated by a
sentinel coefficient $1$ at the crossing count
$s = \texttt{epoch\_end} - \texttt{epoch\_start}$; the sentinel keeps every
committed sequence nonzero (an empty window is the constant $1$), so the
commitment never falls on the identity point, and it pins the sequence's
exact length. Step \texttt{NullifierStep} builds its single-leaf commitment
homomorphically, $[\nfsym]\,\mathcal{G}_0 + \mathcal{G}_1$, with
$\mathcal{G}_1$ carrying the sentinel. Composition is checked as polynomial
identities at a Fiat--Shamir challenge:
\[
C(X) = L(X) + X^{s}\,\bigl(R(X) - 1\bigr)
\qquad \text{(\texttt{UnspentFuse}),}
\]
\[
C(X) = L(X) + X^{s}\,(p - 1) + X^{s+1}\,R(X)
\qquad \text{(\texttt{UnspentEpochFuse}),}
\]
\[
R(X) = E(X) + X^{s}\,\bigl(\texttt{nf\_end} - 1\bigr) + X^{s+1}
\qquad \text{(\texttt{VerifyUnspent}),}
\]
where $p$ is the left half's completing tip \texttt{nf\_end} and $E$ the
\texttt{elapsed} sequence: in each identity the appended term overwrites
the left sentinel and the trailing term re-terminates. The inputs $L$,
$R$, and $p$ are bound by the recursive verification of the input PCDs
\emph{before} the challenge, and each identity is linear in them --- which
is what makes the splice sound. Step \texttt{MergeStamp} fuses two stamps at
equal anchors by requiring, for each of the action-digest and tachygram
multisets, that the merged set polynomial equal the \emph{product} of the
two inputs' (book, \texttt{src/proof-tree.md} @ \texttt{26fdc165}).
\end{construction}

\paragraph{Lineage and spend binding.} Step \texttt{SpendableInit} checks
$\cmsym$ among the creation stamp's tachygrams and roots the chain at
$\texttt{pre\_epoch\_anchor.next\_epoch(epoch)}$, pinning the starting GGM
leaf index to the consensus epoch --- without it, a note spent in its
creation epoch crosses no boundary and the index is a free witness. On
\texttt{present\_nf} the book disagrees with itself: its walkthrough binds
\texttt{present\_nf} to the proven leaf by equality at init --- matching
the step's \texttt{NullifierHeader} input in
Table~\ref{tab:tachyon-design-steps} --- while its reference section calls
it ``unconstrained here'', pinned to a genuine leaf only by the first lift
--- another unresolved internal inconsistency at \texttt{26fdc165}
(\texttt{src/proof-tree.md}). Step \texttt{SpendBind} requires
$\texttt{spendable.cm} = \texttt{note.commitment()}$ --- rejecting a
phantom note reusing $\psi$ under a different value and $\cmsym$ --- and
its \texttt{SpendHeader} is consumed only by \texttt{SpendStamp}, so the
note never propagates. Step \texttt{SpendStamp} composes it with a length-2
\texttt{NullifierHeader} range, witnesses $\texttt{nf\_next}$, binds the
published pair to the boundary leaves by equality, and guards both
nullifiers nonzero (the guard's rationale is quoted at
\S\ref{sec:tachyon-dbu-nf}); deferring this binding and the action
digest here keeps each step within its per-step gate budget. The book's
privacy claim: the note's age never becomes public --- the lineage carries
a single current nullifier, not a polynomial with a consumed offset.

\paragraph{Stamps.} Step \texttt{OutputStamp} is the only stamp-producing
step taking an anchor as direct witness (an output has no prior chain
state); it rejects zero and over-range values, and the wallet typically
anchors outputs at the spends' height so \texttt{MergeStamp} proceeds
without an intervening lift. Step \texttt{StampLift} advances a composed
stamp's anchor toward the tip before publication; validators reconstruct
both set commitments from published bundles, check the proof against them,
and confirm the anchor against the chain. An aggregated stamp has the same
shape as any other and is eligible for further aggregation.

The Orchard contrast: Orchard proves membership on every spend via a
32-layer Sinsemilla Merkle path against a per-block tree anchor (the
\emph{Ironwood Guide}, \S``Proof of ownership of \emph{some} valid note: the
commitment tree''); the no-tree, proven-once replacement is stated at
\S\ref{sec:tachyon-removes}.

\paragraph{Classification.} \emph{Designed but unspecified} throughout:
the nineteen steps exist as circuit-shaped Rust against the mock
(\texttt{stamp/proof/mod.rs}), not as mathematics, and the soundness
section is prose (\S\ref{sec:tachyon-fold};
\S\ref{sec:tachyon-dbu-rollover}). A specification must pin the header
encodings, the per-step relations, the polynomial-commitment scheme and
coefficient generators $\mathcal{G}_i$ (the book's $\Comm$ is abstract),
and the Poseidon instance.

\FloatBarrier

\subsection{Aggregation}\label{sec:tachyon-design-agg}

The lifecycle nomenclature is fixed in
Table~\ref{tab:tachyon-design-agg}: wallets broadcast \emph{autonomes};
aggregators strip and merge stamps into \emph{aggregates}; miners strip
covered transactions into \emph{adjuncts}, and a block carrying an
aggregate must include every covered transaction sans its now-redundant
proof. Aggregates subdivide into \emph{innocent} (no Tachyon actions of
its own, e.g.\ other-pool-only activity) and \emph{based} (proving its own
actions, e.g.\ a coinbase aggregate paying miner rewards through a Tachyon
action). The process is decompress--merge--recompress: select
transactions; deserialize and validate stamps; match or update anchors
(\texttt{StampLift}); merge stamps --- tachygrams as exclusive sets,
proofs recursively; serialize, compress, publish (book,
\texttt{src/aggregation.md} @ \texttt{26fdc165}).

\begin{table}[htbp]
\centering
\small
\begin{tabular}{@{}llll@{}}
\toprule
Term & State & Provenance & Validation \\
\midrule
autonome  & stamped  & wallet     & complete proof with all inputs \\
aggregate & stamped  & aggregator & merged proof needs input from other
  transactions \\
adjunct   & stripped & miner      & input to merged proof from another
  transaction \\
\bottomrule
\end{tabular}
\caption{Bundle lifecycle nomenclature (book, \texttt{src/aggregation.md}
@ \texttt{26fdc165}).}
\label{tab:tachyon-design-agg}
\end{table}

\paragraph{The action-set indicator.} An aggregate's stamp proves its own
plus every covered adjunct's actions but carries no proof header ---
verifiers reconstruct it from covered actions, which before block
inclusion are not at hand. The stamp trailer therefore carries
\texttt{cActionsTachyon}, an author-provided \emph{assistive indicator}
mirroring the action-set commitment the proof attests to, letting an
observer cheaply tell an autonome from an aggregate and judge a collection
of adjuncts complete without verifying the proof. It is authorization
data, not effecting data --- it rides the strippable trailer, contributes
to \texttt{auth\_digest}, and is absent from adjuncts --- and \emph{not} a
soundness mechanism: the proof binds the action set regardless, and a
wrong indicator only harms its author. Function \texttt{Stamp::verify}
uses it as a fast-fail gate before expensive proof verification (book,
\texttt{src/aggregation.md}; code, \texttt{stamp/mod.rs}).

The Orchard contrast: no analogue exists --- an Orchard bundle's proof is
per-transaction, never stripped, never merged across transactions (the
\emph{Ironwood Guide}, \S``What a node verifies, without ever observing the
note'').

\paragraph{Classification.} \emph{Designed but unspecified}, shading into
\emph{open problem} inside the chapter's own file: the tachygram set-union
validation algorithm and the aggregation limits are in-file TODOs, and the
closing comment lists undecided economics --- aggregate size limited by
commitment size before block size (``on the order of thousands of
actions''), single-aggregate merging serial though multiple aggregates
parallelize, ser/de expensive, and the optimal-aggregate,
miner-integrate-vs-p2p, and selection-order questions open
(\texttt{src/aggregation.md}, comments). The ``Block Layout'' section the
transaction-identifiers chapter cites does not exist in
\texttt{aggregation.md} at this commit --- a dangling link. A
specification must pin the union-validation algorithm, the limits, and
the gossip protocol.

\subsection{Authorization}\label{sec:tachyon-design-auth}

A bundle carries three authorization layers: per-action signatures binding
each tachyaction to its tachygram, value commitments hiding individual
values while preserving their algebraic sum, and a binding signature
proving the declared balance (book, \texttt{src/authorization.md} @
\texttt{26fdc165}).

\begin{construction}[Deterministic action randomizer]
\label{def:tachyon-design-alpha}
The planner selects arbitrary per-action entropy $\theta$; the custody
device receives each note and $\theta$ to independently confirm the
planning work. The randomizer is deterministic from $(\theta, \cmsym)$
under per-effect personalizations:
\[
\alpha_{\mathrm{spend}} =
  \text{BLAKE2b-512}_{\texttt{Tachyon-Spend}}(\theta \,\|\, \cmsym),
\qquad
\alpha_{\mathrm{output}} =
  \text{BLAKE2b-512}_{\texttt{Tachyon-Output}}(\theta \,\|\, \cmsym).
\]
Every action signs with a unique $\rsk$, published as
$\rk = [\rsk]\,\mathcal{G}$. For a spend,
\[
\rsk = \ask + \alpha, \qquad
\rk = [\ask + \alpha]\,\mathcal{G} = \aksym + [\alpha]\,\mathcal{G},
\]
so the planning device derives $\rk$ from the validating $\aksym$ without
ever holding $\ask$, and custody confirms $\rk$ before signing. For an
output, $\rsk = \alpha$ and $\rk = [\alpha]\,\mathcal{G}$: no authority is
needed. Because $\alpha$ derives from $\theta \,\|\, \cmsym$, the key $\rk$
is itself a commitment to $\cmsym$, and the signature transitively binds
each action to its tachygram without the tachygram appearing in the action.
\end{construction}

\paragraph{Bundle commitment and sighash.} The bundle's effect digest
accumulates per-action digests into a set polynomial:
\[
d_i = \mathsf{Poseidon}_{\texttt{Tachyon-ActionDg}}
  (\cvsym_i \,\|\, \rk_i), \qquad
\mathsf{action\_acc} = \Comm\Bigl(\prod_i \bigl(X - d_i\bigr)\Bigr),
\]
hashed as $\text{BLAKE2b-512}_{\texttt{ZTxIdTachyonHash}}
(\mathsf{action\_acc} \,\|\, \mathsf{value\_balance})$. The same
action-digest accumulation is bound to the stamp proof; the stamp itself
is excluded from the bundle commitment because it strips during
aggregation. All signatures --- action and binding --- sign one
transaction-wide sighash, computed at the transaction layer from every
pool's bundle commitment and passed to the Tachyon layer as opaque bytes.
Consensus recomputes the action digests from the visible actions and
checks them against both the sighash (via the bundle commitment) and the
stamp (via the polynomial commitment in the PCD header); a modified action
breaks both checks.

\paragraph{Value balance.} The value commitment and binding algebra are
Orchard's, deliberately: with $v$ signed (positive spends, negative
outputs) and $\rcv$ random in $\F_{p_{\Vesta}}$,
\[
\cvsym = [v]\,\mathcal{V} + [\rcv]\,\mathcal{R}, \qquad
\mathsf{bvk} = \sum_i \cvsym_i - [\mathsf{value\_balance}]\,\mathcal{V}
 = \Bigl[\sum_i v_i - \mathsf{value\_balance}\Bigr]\,\mathcal{V}
   + [\mathsf{bsk}]\,\mathcal{R},
\]
where $\mathsf{bsk} = \sum_i \rcv_i$: the planner signs the sighash with
$\mathsf{bsk}$ directly, without custody assistance, and validators'
reconstructed $\mathsf{bvk}$ equals $[\mathsf{bsk}]\,\mathcal{R}$ exactly
when $\sum_i v_i = \mathsf{value\_balance}$. The generators are
\emph{shared with Orchard}, from the domain \texttt{z.cash:Orchard-cv}.
Two footnotes flag churn: the $\rcv$ derivation ``may be revised in the
future to incorporate a hash of the note commitment'', and the binding
signature uses \texttt{reddsa::orchard::Binding}, hardcoding $\mathcal{R}$
--- ``We should consider defining a unique personalization.''

The Orchard contrast: the value machinery is identical by construction ---
same Pedersen $\cvsym$, same $\mathsf{bvk}$ equation, literally the same
generators (the \emph{Ironwood Guide}, \S``Conservation of value: value
commitments and the binding signature'') --- while the randomizer differs
in kind: Orchard draws $\alpha$ uniformly at random per Action with no
planner/custody split (\emph{Ironwood Guide}, \S``Authorisation of the
spend: RedPallas''), and every Orchard Action is a paired
spend-plus-output carrying one \textsf{SpendAuthSig} (\emph{Orchard
Guide}, \S``A single shielded payment, end to end''; \S``What a node
verifies, without ever observing the note'').

\paragraph{Classification.} This is the \emph{specified}-as-code side of
the README split: the randomizer, re-randomized keys, value commitments,
and both signature layers are implemented and tested (\texttt{entropy.rs},
\texttt{keys/private.rs}, \texttt{value.rs}, \texttt{reddsa.rs}). Unpinned
by any artifact: the sighash construction (opaque bytes at this layer),
the $\rcv$ derivation, the generator personalization, the $\Comm$ scheme
behind $\mathsf{action\_acc}$, and --- flagged in-file --- whether
$\alpha$ moves to a curve-native derivation
(\texttt{src/domain-separators.md}, comment).

\subsection{Transaction identifiers and domain separators}
\label{sec:tachyon-design-txid}

A bundle's authorization form changes through aggregation --- stamping,
merging, and stripping produce bit-different authorizations of the same
effecting data --- and \texttt{wtxid} $=$ \texttt{txid} $\|$
\texttt{auth\_digest} (ZIP~239) must distinguish the forms. Tachyon routes
every mutable part through \texttt{auth\_digest}: \texttt{txid} commits to
effecting data only, $\mathsf{action\_acc} \,\|\, \mathsf{value\_balance}$,
so stripping, merging, and re-stamping leave it unchanged, while the
authorization side commits to the signatures plus the stamp trailer:
\[
\begin{split}
\mathsf{auth\_digest\_contribution} &=
\text{BLAKE2b-256}_{\texttt{ZTxAuthTachyHash}}\bigl(
  \mathsf{vActionSigs} \| \mathsf{bindingSig} \| T\bigr), \\
T &= \begin{cases}
    \texttt{cActionsTachyon}_{32} \| \texttt{anchor}_{32} \|
      \texttt{vTachygrams} \| \texttt{proof} & \text{if stamped} \\
    \texttt{stampWtxid}_{64} & \text{if stripped}
  \end{cases}
\end{split}
\]
(book, \texttt{src/transaction-identifiers.md} @ \texttt{26fdc165}). Each
physical form thus yields a distinct \texttt{wtxid}. The personalization
\texttt{ZTxAuthTachyHash} is declared ``a placeholder until a Tachyon-ZIP
amendment to ZIP-244 fixes it.'' An adjunct's covering-aggregate reference
is a \texttt{wtxid}, assigned by the miner at block assembly; the covering
aggregate must itself be top-level in the block --- never stripped, never
further aggregated --- so the reference is stable.

\paragraph{A book-internal (and book--code) discrepancy.} The
transaction-identifiers formula writes BLAKE2b-\emph{256} for the
\texttt{auth\_digest} contribution --- matching ZIP~244's digests ---
while the book's own domain-separator table
(Table~\ref{tab:tachyon-design-domains}) and the crate
(\texttt{hash\_length(64)}; \texttt{digest/blake2b.rs}) both say
BLAKE2b-\emph{512}: the hash width is an unresolved encoding at
\texttt{26fdc165}. The code also fixes encodings the book leaves loose:
the bundle commitment hashes the action commitment's complete coordinates
$x \,\|\, y$ then \texttt{value\_balance} as a little-endian i64, and a
transaction with no Tachyon bundle contributes the personalized hash of
the \emph{empty} string --- documented as distinct both from a stripped
bundle and from a bundle with no actions and zero balance
(\texttt{blake2b.rs}, \texttt{bundle/mod.rs}).

The Orchard contrast: the construction mirrors the ZIP-244
personalized-BLAKE2b pattern of \texttt{ZTxIdOrchardHash} (the
\emph{Ironwood Guide}, \S``Conservation of value: value commitments and the
binding signature'', paragraph ``The transaction digest (ZIP-244)''); the
novelty is that the entire stamp rides \texttt{auth\_digest}, keeping
\texttt{txid} invariant under aggregation.

\paragraph{Domain separators.} Table~\ref{tab:tachyon-design-domains} is
the book's complete table --- its most specification-like artifact,
byte-exact strings each realized as a code constant. The crate fixes the
in-circuit realization the book omits: every Poseidon domain tag is a
16-byte ASCII string absorbed as the \emph{first} sponge input, embedded
as $\F_{p_{\Pallas}}$ via
\texttt{Fp::from\_u128(u128::from\_le\_bytes(tag))}, with arities: action
digest 5 (domain, $\cvsym$ and $\rk$ as full coordinate pairs), payment
key 3, note commitment 5, master key 3, prefix step 3, nullifier 2, anchor
stamp step 6, anchor empty step 2, anchor epoch step 3
(\texttt{digest/poseidon.rs}) --- all through the mock sponge.

\begin{table}[htbp]
\centering
\footnotesize
\begin{tabular}{@{}lll@{}}
\toprule
Purpose & Hash & Value \\
\midrule
Spend $\alpha$          & BLAKE2b-512 & \texttt{Tachyon-Spend} \\
Output $\alpha$         & BLAKE2b-512 & \texttt{Tachyon-Output} \\
Bundle commitment       & BLAKE2b-512 & \texttt{ZTxIdTachyonHash} \\
Bundle auth digest      & BLAKE2b-512 & \texttt{ZTxAuthTachyHash} \\
PRF expand ($\mathtt{0x21}$ $\ask$, $\mathtt{0x22}$ $\nksym$)
                        & BLAKE2b-512 & \texttt{Zcash\_ExpandSeed} \\
\midrule
Nullifier prefix key    & Poseidon & \texttt{Tachyon-NfPrefix} \\
Nullifier derivation    & Poseidon & \texttt{Tachyon-NfDerive} \\
Note commitment         & Poseidon & \texttt{Tachyon-CmDerive} \\
Action digest           & Poseidon & \texttt{Tachyon-ActionDg} \\
Payment key derivation  & Poseidon & \texttt{Tachyon-PkDerive} \\
Anchor stamp step       & Poseidon & \texttt{Tachyon-StampFld} \\
Anchor empty step       & Poseidon & \texttt{Tachyon-EmptyBlk} \\
Anchor epoch step       & Poseidon & \texttt{Tachyon-EpochStp} \\
\midrule
Value commitment        & hash-to-curve & \texttt{z.cash:Orchard-cv} \\
\bottomrule
\end{tabular}
\caption{Complete domain-separator table (book,
\texttt{src/domain-separators.md} @ \texttt{26fdc165}). The auth-digest
hash width is the discrepancy noted in the text: the book's
transaction-identifiers formula says BLAKE2b-256; this table and the
crate say BLAKE2b-512.}
\label{tab:tachyon-design-domains}
\end{table}

\paragraph{Classification.} \emph{Designed but unspecified}, at the highest
concreteness the corpus reaches: byte-exact strings with code constants,
yet two entries are self-declared placeholders ---
\texttt{ZTxAuthTachyHash}, and the value-commitment generator borrowed
from Orchard's \texttt{z.cash:Orchard-cv} --- and one is internally
inconsistent (the 256/512 width, above).
A specification must fix the hash widths, the personalization
bytes, and above all the Poseidon instance the eight domains parameterize
--- parameters, rate, and capacity exist nowhere, and the crate computes
every domain through a mock.

\medskip
\noindent
Taken together, the book at \texttt{26fdc165} fixes an object model at a
concreteness no other Tachyon artifact approaches, and still leaves every
load-bearing parameter symbolic: the GGM shape and epoch length
(\S\ref{sec:tachyon-dbu-epoch}), the proof size, the Poseidon instance,
the commitment scheme and its generators, the sighash, and the
note-commitment scheme itself. The sections that follow classify these
gaps severally; the open problems are the design's own
(\S\ref{sec:tachyon-open}).

\section{What is specified today}\label{sec:tachyon-specified}

Every claim in this volume is classified \emph{specified}, \emph{designed
but unspecified}, or \emph{open problem}. This section is the
specified bin, and it is exhaustive: a statement appears here only if a
verifiable artifact backs it --- a ZIP with a recorded status in
\texttt{zcash/zips}, or source code read at a pinned commit. Our ground truth
is \texttt{zcash/zips} at commit \texttt{c6b70358} (2026-07-05),
\texttt{tachyon-zcash/ragu} at \texttt{830bbcda} (2026-07-05), and
\texttt{tachyon-zcash/tachyon} at \texttt{26fdc165} (2026-07-07); repository
metadata (releases, package registries, continuous integration) was queried
on 2026-07-08. The two lower bins follow in \S\ref{sec:tachyon-dbu}
(designed but unspecified) and \S\ref{sec:tachyon-open} (open problems).

The headline fact is negative. No component of the Tachyon core --- the
shielded protocol, the bundle format, the accumulator, the aggregator,
oblivious synchronization --- is specified anywhere: upstream
\texttt{zcash/zips} contains no Tachyon ZIP or pull request, and the five
core ZIP stubs in the project's own tracker are empty files
(\texttt{book/src/zips/} @ \texttt{26fdc165}; tachyon-zcash/tachyon
issue \#114). What the specified bin actually contains is periphery: one
Active payment-request substrate, four groundwork drafts in varying states
of health, and the code state of the proof layer and the protocol crate. We
inventory each in turn.

\subsection{ZIP 321 as substrate, and the four groundwork drafts}
\label{sec:tachyon-groundwork}

All ZIP facts below are read from \texttt{zcash/zips} @ \texttt{c6b70358};
line numbers refer to the files at that commit.

\paragraph{ZIP 321 --- Payment Request URIs (Status: Active).}
This is the only entry at deployed rigor. It defines the \texttt{zcash:} URI
scheme by which a payer is asked for payment: a grammar of
\texttt{address}, \texttt{amount}, \texttt{memo} (base64url),
\texttt{label}, and \texttt{message} parameters, with numeric parameter
indices for multi-recipient requests (\texttt{zip-0321.rst}). Its role here
is substrate: the pre-pivot Tachyon design removes per-note key
diversification, and the design book assigns per-payment unlinkability to
``out-of-band payment protocols (ZIP 321 payment requests, ZIP 324 URI
encapsulated payments)'' (tachyon \texttt{book/src/keys.md} @
\texttt{26fdc165}).

\paragraph{ZIP 324 --- URI-Encapsulated Payments (Status: Draft).}
Created 2019-07-17, the draft predates Tachyon by six years and mentions it
nowhere. It specifies sending a payment through any secure messaging channel
as a URI encoding the secret spending key of an ephemeral, freshly funded
address; anyone who learns the URI finalizes the payment by sweeping the
encapsulated funds into their own wallet (\texttt{zip-0324.rst}, Abstract).
Three limits bound what it pins down. First, the construction is
Sapling-only: the encapsulated note is a Sapling note, keys derive under the
path \texttt{m\_Sapling/324'/\ldots}, and Orchard appears zero times.
Second, core details are unpinned: the master-key derivation
$\mathsf{rseed} = \mathsf{PRFexpand}_{\mathsf{sk}_m}(\cdot)$ carries the
literal placeholder \texttt{[0xFIXME]} as its domain separator (line 148),
the note-identification section is marked ``out of date \ldots{} To be
revised'' (line 384), and the draft closes with a ``Publication Blockers''
section (line 458). Third, the hard part is out of scope by declaration:
``It is outside the scope of this proposal to establish a secure
communication channel for transmission of URI-Encapsulated Payments''
(line 107) --- the unlinkability Tachyon defers to this document is thus
deferred again, to the messaging layer.

\paragraph{ZIP 231 --- Memo Bundles (Status: Draft).}
The cryptographic core is at full ZIP rigor. Memos leave the note plaintext
and move into a per-transaction \emph{memo bundle}: a sequence of $272$-byte
encrypted chunks capped at
$\mathsf{memo\_chunk\_limit} \times 256 = 16384$ bytes of plaintext
($16$\,KiB) per transaction (lines 209--216); each chunk sealed with
ChaCha20-Poly1305 in a variant of the STREAM construction, the final chunk
distinguished by a nonce suffix (lines 264--274); per-output $32$-byte memo
keys $\mathsf{K^{memo}}$ with two reserved values --- all-$\mathtt{0x00}$
for publicly decryptable, all-$\mathtt{0xFF}$ for no memo (lines 222--231);
chunks of all memos interleaved by an order-preserving shuffle that MUST
reveal to a recipient only the size of their own memo and the total count of
chunks they cannot decrypt (lines 282--297); a two-pass decryption
algorithm, one pass under counter nonces and one under the final-chunk nonce
(lines 323--351); independent prunability via $\mathtt{fAllPruned}$, with
relay of pruned transactions forbidden (lines 391--430); and a fee treatment
under ZIP 317 Revision 2, two free chunks whenever the transaction has
shielded outputs (line 422\,ff.). The gaps are equally explicit: the sighash
commitment is ``TODO: finish this as a modification to ZIP 248'' (line 417);
the lead byte is the unassigned placeholder \texttt{\{\{MBLEADBYTE\}\}}
(lines 490, 588); and the carrier is orphaned --- the transaction fields
that would hold a memo bundle were defined by the withdrawn ZIP 230, and
their designated new home, ZIP 248, is unmerged.

\paragraph{ZIP 230 --- Version 6 Transaction Format (Status: Withdrawn).}
The draft carried the byte-precise OrchardZSA v6 wire format: Action Groups,
the ZSA issuance bundle, the memo-bundle carrier fields, an explicit
\texttt{fee} field, \texttt{zip233Amount}, and per-transparent-input
\texttt{vSighashInfo} (\texttt{zip-0230.rst}, field tables). Its withdrawal
warning states that ``Transaction version number 6 is now defined by
ZIP 248'' (lines 39--40) --- already stale, since ZIP 229 (Status: Draft,
created 2026-06-13) now defines v6 for the NU6.3 pool and declares as
non-requirements exactly the withdrawn machinery: ZSAs,
\texttt{zip233Amount}, the explicit \texttt{fee} field, per-input sighash
information, and ZIP 248's ``other extensibility affordances''
(\texttt{zip-0229.md}, Non-requirements). The ZIP 230 byte format survives
only as reference text with no deployment path.

\paragraph{ZIP 246 --- Digests for the Version 6 Transaction Format
(Status: Withdrawn).}
Created 2025-02-12 and withdrawn 2026-06-24, the draft defined v6
transaction identifiers, authorizing-data commitments, and signature
digests; its withdrawal note reassigns all three to ZIP 229 (lines 39--41).
Its one surviving idea is sighash-algorithm versioning: a closed
per-transaction-version set of sighash algorithms, selected via
$\mathtt{sighashInfo} = [\mathtt{sighashVersion}] \,\|\,
\mathtt{associatedData}$ (lines 106--123). The concept now lives only in
withdrawn text and in the unmerged ZIP 248 pull request (zcash/zips\#1156,
idle since 2026-05-26); ZIP 229 itself hashes in the plain ZIP-244 style.

\medskip
\noindent
The pattern across the four drafts is uniform. Each piece of groundwork
either lacks its carrier (ZIP 231), was withdrawn in favor of a narrower
NU6.3 format (ZIPs 230 and 246), or defers its hardest component to another
layer (ZIP 324). The single live v6 specification today is ZIP 229; the
extensible-format home that ZIPs 230, 231, and 246 all defer to, ZIP 248,
exists only as an open pull request.

\subsection{Ragu at commit \texttt{830bbcda}: the proof layer as code}
\label{sec:tachyon-ragu}

Ragu is the proof-carrying-data (PCD) framework on which Tachyon's
wallet-state proofs are to run. No protocol specification exists: every
protocol chapter of the Ragu book --- assumptions, arithmetization, NARK,
split accumulation, registry polynomial, recursion, analysis --- is marked
\texttt{<!-- todo -->} in the book index (\texttt{book/src/SUMMARY.md}).
At this commit the code is therefore the only specification of what Ragu
does, and we read the code.

\paragraph{What the pipeline implements.}
The repository is a \texttt{no\_std} Rust workspace (version 0.0.0, edition
2024, MSRV 1.90) implementing a modified version of the Halo recursive-SNARK
construction, with no trusted setup (\texttt{README.md},
\texttt{Cargo.toml}). The Pasta cycle is the only implemented curve cycle:
crate \texttt{ragu\_pasta} provides the sole \texttt{Cycle} implementation,
with the Pallas base field as the circuit field
(\texttt{crates/ragu\_pasta/src/lib.rs}). A working PCD pipeline exists end
to end: an application registers its steps through
\texttt{ApplicationBuilder::register} and finalizes; seeding creates a leaf;
\texttt{fuse} merges two proof-carrying nodes into one through an
eleven-stage prover pipeline (files \texttt{\_01\_application.rs} through
\texttt{\_11\_circuits.rs} under \texttt{crates/ragu\_pcd/src/fuse/});
\texttt{rerandomize} and \texttt{verify} complete the interface. Claims are
tracked in two registries, native (Pallas-field) and nested (Vesta-field,
CycleFold-style). An integration test builds and verifies a multi-node PCD
tree (\texttt{crates/ragu\_pcd/tests/nontrivial.rs}).

\paragraph{Verification is linear-time; no opening argument exists.}
Proofs are not succinct. Type \texttt{Proof} carries the full sparse
polynomials of the argument --- fifteen native polynomials plus the bridge
polynomials, some twenty in all (\texttt{ragu\_pcd/src/proof/mod.rs}).
Function \texttt{verify()} checks a proof by sampling challenges, rebuilding
the revdot claim polynomials, and evaluating the carried polynomials
directly (\texttt{ragu\_pcd/src/verify.rs}); an in-code note records
that the native $c$-claim is tautological in the verifier and meaningful
only inside the circuit. No inner-product-argument opening, decider, or
proof-compression step exists anywhere in the workspace. Verification
therefore runs in time linear in the proof, and the succinct final verifier
Tachyon needs is future work, not present code.

\paragraph{Soundness placeholders.}
Three placeholders sit on the main branch, each self-documented. The
Fiat--Shamir domain-separation tag is the literal
\texttt{RAGU\_TAG: \&[u8] = b"FIXME"}, beneath the comment ``choose a
permanent domain separation tag before release''
(\texttt{ragu\_pcd/src/lib.rs}, lines 50--51). Registry binding ---
the mechanism that replaces per-circuit verification keys --- runs six hash
iterations from the placeholder ``nothing-up-my-sleeve'' challenges $w = 2$,
$x = 3$, $y = 5$, under an in-code \texttt{FIXME(security)} stating that six
iterations ``is insufficient to fully bind the registry polynomial''
(\texttt{ragu\_circuits/src/registry.rs}, lines 673--692; issues
\#78, \#316). The verifier, finally, omits three checks recorded in an
in-code TODO: \texttt{registry\_wx0\_poly}, \texttt{registry\_wx1\_poly},
and \texttt{registry\_wy\_poly} (\texttt{verify.rs}, lines 122--124).
Beyond soundness, zero knowledge is opt-in rather than default: the
documentation of \texttt{fuse()} warns that accepting an RNG ``is not an
indication that the resulting proof-carrying data is zero-knowledge; that
must be ensured by performing \texttt{Application::rerandomize} at a later
point'' (\texttt{fuse/mod.rs}), and \texttt{rerandomize} itself seeds a
trivial proof under a comment reading ``this is a temporary hack''
(\texttt{lib.rs}, line 262).

\paragraph{Assurance machinery.}
The engineering hygiene around this pre-release core is heavy and, at the
pinned commit, green. The workspace carries 445 \texttt{\#[test]} functions
and twelve property-testing suites; line coverage stood at 93.8\,\%
(Codecov snapshot, 2026-03-08). Twenty libFuzzer targets --- witness
generation and cheating, revdot and algebraic identities, Poseidon
differentials, metamorphic driver behavior, verifier rejection, among
others --- run on a twice-weekly cron
(\texttt{qa/fuzz/fuzz\_targets/},
\texttt{.github/workflows/fuzz-cron.yml}: Sundays and Wednesdays). A Lean~4
formalization covers twenty-seven gadget-circuit files, and CI checks that
the Lean sources extracted from the Rust gadgets are current
(\texttt{qa/lean/Ragu/Circuits/}, \texttt{.github/workflows/fv.yml}).
Supply-chain review runs under \texttt{cargo-vet}
(\texttt{qa/supply-chain/}) and workflow linting under \texttt{zizmor}. All
continuous-integration workflows were green at \texttt{830bbcda} (GitHub
Actions, queried 2026-07-08).

\paragraph{What is absent.}
No security audit exists: the README warns, verbatim, ``Ragu is under heavy
development and has not undergone auditing. Do not use this software in
production.'' No benchmark has been published: a bench harness runs in CI
(\texttt{.github/workflows/bench.yml}) but publishes no numbers, and the
only performance figures anywhere in the project are the Ragu book's
self-labelled rough estimates --- which belong to the
designed-but-unspecified bin, not to this one. No release exists: the
workspace version is 0.0.0, the repository has zero GitHub releases, and no
\texttt{ragu} crate is on crates.io (registry queries, 2026-07-08).

\subsection{The \texttt{zcash\_tachyon} crate: integration against a mock}
\label{sec:tachyon-crate}

Crate \texttt{zcash\_tachyon} (tachyon-zcash/tachyon @ \texttt{26fdc165})
implements and tests the conventional cryptographic plumbing: key derivation
via $\mathsf{PRFexpand}$ with sign normalization, Pedersen value commitments
over Pallas sharing Orchard's generators, randomized RedPallas
spend-authorization signatures, binding signatures, and bundle construction
(\texttt{README.md}). Everything proof-shaped is stubbed. The crate consumes
Ragu exclusively through its \texttt{mock} feature, pinning the dependency
to a contributor fork (\texttt{github.com/turbocrime/ragu}, rev
\texttt{98a15d32}) with features \texttt{mock} and \texttt{legacy-deps}
(\texttt{crates/tachyon/Cargo.toml}). The mock is, in its own words, an
``API-level mock of \texttt{ragu\_pcd}'' that ``mirrors the shape of the
real \texttt{ragu\_pcd} API so downstream consumers (e.g.\ Zebra) can
integrate against it ahead of the real implementation'' (ragu
\texttt{src/mock/mod.rs} @ \texttt{830bbcda}); it performs no cryptography.
Against this interface the crate registers its nineteen proof steps
(\texttt{crates/tachyon/src/stamp/proof/mod.rs}). The README's own stub
list is exact: note commitments, nullifier derivation (the GGM tree PRF),
proof creation, merge, and verification, and stamp
compression/decompression all await ``Ragu PCD and Poseidon integration''.
The shape of the integration is therefore specified by code; its
cryptographic content is not.

\medskip
\noindent
Taken together, the specified bin holds one Active URI substrate; one
out-of-band payment draft that is Sapling-only with its core derivation
unpinned; one memo-bundle draft whose cryptography is complete but whose
carrier is orphaned; two withdrawn format ZIPs surviving as reference text;
a proof framework whose pipeline runs but whose verifier is linear-time,
incomplete, and Fiat--Shamir-tagged \texttt{b"FIXME"}; and a protocol crate
integrating against a no-crypto mock. Every load-bearing Tachyon design ---
evolving nullifiers, the proof tree, aggregation, oblivious
synchronization --- lives outside this bin, in the sections that follow.

\section{Designed but unspecified, and open problems}\label{sec:tachyon-unspecified}

This section carries the two lower bins of the classification of
Definition~\ref{def:three-bins}: constructions that exist in an informal
source but in no artifact at ZIP or protocol-specification rigor, and
problems the designers themselves acknowledge as unsolved. Nothing here is
normative, and the global fact that forces both bins is established in
\S\ref{sec:tachyon-specified}: upstream \textsf{zcash/zips} (commit
\texttt{c6b70358}, 2026-07-05) contains zero Tachyon ZIPs or pull requests,
and the five core ZIP pages in the project's own design book --- shielded
protocol, bundle, accumulator, aggregator, and oblivious syncing service ---
are eleven-line skeletons with empty \emph{Dependencies}, \emph{Design
Considerations}, and \emph{ZIP Draft} headings (\texttt{book/src/zips/} @
\texttt{26fdc165}, 2026-07-07).

For every entry we state three things: \emph{what exists}, with the source
and commit (source pinning as in \S\ref{sec:tachyon-sources}; the mechanisms
themselves are constructed in \S\ref{sec:tachyon-architecture}); \emph{what
a specification must pin down} before the entry can move to the specified
bin (for open problems, what a resolution must establish); and \emph{what
would falsify the design} --- the observation that would not merely delay
the entry but refute it.

\subsection{Designed but unspecified}\label{sec:tachyon-dbu}

\subsubsection{Nullifier evolution and re-randomization for oblivious
synchronisation}\label{sec:tachyon-dbu-nf}

\paragraph{What exists.} The ePrint note introduces evolving nullifiers only
as a sketch: ``As a concrete example, assume the nullifier for preventing
the double-spend of a note is modified from $\eta \leftarrow
\mathrm{PRF}(k_{\mathrm{note}}, \rho)$ to $\eta_e \leftarrow
\mathrm{PRF}(k_{\mathrm{note}}, \rho \| e)$'' with $e$ ``a monotonic counter
known as an epoch identifier'' (ePrint 2025/2031, \S1.3). The concrete
derivation exists only in the design book
(\texttt{book/src/nullifiers.md} @ \texttt{26fdc165}): a
Goldreich--Goldwasser--Micali tree of depth $D$ and arity $A$ rooted at the
per-note master key
\[
\mathsf{mk} = \mathsf{Poseidon}_{\texttt{Tachyon-NfPrefix}}(\psi,\, \nksym),
\]
descended one Poseidon call per base-$A$ chunk of the epoch index, with the
epoch-$e$ nullifier one further call under a separate domain,
$\nfsym = \mathsf{Poseidon}_{\texttt{Tachyon-NfDerive}}\!\bigl(
\mathsf{KDF}^{\mathsf{climb}}_\psi(e, D)\bigr)$. The note commitment
$\cmsym = \mathsf{Poseidon}_{\texttt{Tachyon-CmDerive}}(\rcm, \mathsf{pk},
v, \psi)$ digests $\psi$ and (through $\mathsf{pk}$) $\nksym$, so one
commitment freezes the note's entire nullifier sequence. Delegation hands a
service a window of future nullifiers committed on coefficient generators
with a sentinel, $\delta = \sum_i [\Delta_{e+i}]\,\mathcal{G}_i +
\mathcal{G}_d$, re-based so any window stands alone. None of this is
implemented: the \textsf{zcash\_tachyon} crate lists ``Note commitments''
and ``Nullifier derivation (GGM tree PRF)'' under ``Stubbed (pending Ragu
PCD and Poseidon integration)'' (\texttt{README.md} @ \texttt{26fdc165};
\S\ref{sec:tachyon-crate}), and \texttt{crates/tachyon/src/note.rs} calls
the commitment scheme ``TBD''. The material is single-committer and
churning: alternatives remain open in the tracker for the derivation itself
(issues \#86, \#139) and for the absence-proof ranges it feeds (\#87, \#97,
\#98; all open, fetched 2026-07-08).

\paragraph{What a specification must pin down.} The depth $D$ and arity $A$
(symbolic in the book; the crate pins working values,
\S\ref{sec:tachyon-design-ggm}); the Poseidon instances,
their parameters, and the byte-exact domain tags; the encoding of the epoch
index into chunks; the uniqueness rule for $\psi$ --- the book states that
two notes sharing $\psi$ share $\mathsf{mk}$ and hence the same nullifier
sequence, ``so spending one publishes the other's nullifiers''
(\texttt{nullifiers.md}), yet no rule assigns responsibility for preventing
reuse; the generators $\mathcal{G}_i$ and the window-commitment format; and
the nonzero guards on published nullifiers, which the book motivates only in
prose (a zero nullifier ``would collide with the note's own $\cmsym$ in the
tachygram scan'', \texttt{proof-tree.md}).

\paragraph{What would falsify the design.} A distinguisher on the per-epoch
values: obliviousness rests entirely on nullifiers at distinct epochs being
unlinkable, so a practical distinguisher for
$\mathsf{Poseidon}$-as-PRF in this mode --- linking $\nfsym_e$ to
$\nfsym_{e'}$ without $\mathsf{mk}$ --- collapses the construction. Likewise
a derivation-chain break: a valid absence proof over values that are not the
note's genuine leaves would let a spender evade the double-spend check the
whole system replaces the nullifier set with.

\subsubsection{The PIR detection design}\label{sec:tachyon-dbu-pir}

\paragraph{What exists.} One sentence. Issue \#121 (2026-05-26, verified via
the GitHub API 2026-07-08) records the pivot in full: ``Since we're using
PIR to make repeated payments flows without OOB, there are stale parts of
the book that will need to change. We're retaining in-band secret
distribution for recover from seed flows.'' The same issue flags the book's
key-hierarchy chapter as stale, and the chapter carries the marker
\texttt{<!--\ todo: this is significantly outdated -->}
(\texttt{book/src/keys.md} @ \texttt{26fdc165}). Everything else written
about payment detection predates this pivot and describes the abandoned
fully-out-of-band design (blog ``Tachyaction at a Distance'', 2025-05-15):
removal of key diversification, viewing keys, and payment addresses, with
unlinkability deferred to ZIP-321 payment requests and ZIP-324 encapsulated
payments --- the latter itself a 2019 draft whose rseed domain separator is
the literal \texttt{0xFIXME} and which declares the transport channel out of
scope (\textsf{zips} @ \texttt{c6b70358}). The project roadmap describes the
PIR-plus-post-quantum-key-exchange payment protocol as developed
independently in collaboration with the team; no document of any rigor
describes it (tachyon.z.cash/roadmap, dateless, fetched 2026-07-08).

\paragraph{What a specification must pin down.} The PIR family
(single-server computational or multi-server information-theoretic, with
the corresponding non-collusion assumption); the database the queries run
against, who maintains it, and how a detection index is derived from a
payment; the query and server cost model as the database grows; the
post-quantum key exchange (named nowhere); and the retained in-band channel
--- its ciphertext format, its interaction with the removal of viewing keys,
and which flows are entitled to it. The design currently consists of a
sentence naming a primitive; every one of these is unwritten.

\paragraph{What would falsify the design.} Cost: single-server PIR touches
the whole database per query, so a database that grows with chain history
can make detection more expensive than the trial decryption Tachyon set out
to eliminate. Privacy: query-pattern or timing leakage at the PIR server
re-introduces exactly the correlation channel of
\S\ref{sec:tachyon-open-timing}. Either failure, demonstrated at realistic
scale, refutes the pivot rather than the parameters.

\subsubsection{The sync-service trust model}\label{sec:tachyon-dbu-oss}

\paragraph{What exists.} Role prose. The ePrint note asserts services are
fully oblivious to their clients' transaction details and keep ephemeral
per-client state; the book's role table is the most precise statement: the
sync service runs \textsf{UnspentSeed}, \textsf{EmptyBlockUnspentSeed},
\textsf{UnspentFuse}, and \textsf{UnspentEpochFuse} over nullifier values
the wallet shared, and ``never sees a note, \texttt{cm}, \texttt{psi}, or
\texttt{mk}'' (\texttt{book/src/proof-tree.md} @ \texttt{26fdc165}). The
security requirement is stated by the designers in the negative: ``it is
essential that the syncing service be incapable of distinguishing between
transactions they assisted with and ones they did not'' (ePrint 2025/2031,
\S1.3). The mechanism offered is epoch unlinkability plus proof
re-randomization, the latter by citation: the problem ``can be generally
addressed by re-randomizing the proof, as many recursive proof systems
permit, such as those involving accumulation schemes'' (ibid., citing
B\"unz--Chiesa--Mishra--Spooner, ePrint 2020/499). On the implementation
side, \textsf{ragu} does expose \texttt{Application::rerandomize}
(\S\ref{sec:tachyon-ragu}), but its own API documentation warns that fused
proof-carrying data is not zero-knowledge until re-randomization is
performed, and the re-randomization path rests on an in-code ``temporary
hack'' (\texttt{crates/ragu\_pcd/src/fuse/mod.rs}, \texttt{lib.rs} @
\texttt{830bbcda}). The corresponding ZIP stub is empty
(\texttt{book/src/zips/oss.md}).

\paragraph{What a specification must pin down.} A definition of
``oblivious'' as a game --- what the adversarial service observes (requests,
timing, shared windows), what it must not distinguish, and against which
collusion (service alone, service plus network observer, several services
comparing notes); what state a service may retain across sessions,
``ephemeral'' being currently a one-word promise; the client--service
protocol itself, including handoff between services, which no document
describes; and whether re-randomization is mandatory per exchange and on
which side. No security definition of any of this exists in the corpus.

\paragraph{What would falsify the design.} The designers' own criterion,
inverted: a service strategy that distinguishes assisted from non-assisted
transactions. The timing-correlation channel of
\S\ref{sec:tachyon-open-timing} is the known candidate, and the designers
state plainly that its success would be fatal.

\subsubsection{Epoch parameterization}\label{sec:tachyon-dbu-epoch}

\paragraph{What exists.} An open issue. The epoch period --- the single
parameter the pruning model, the delegation window, the dual-nullifier rule,
and the data-availability pressure all hinge on --- is untuned: issue \#79
(``Tune epoch period'', 2026-05-16, verified via the GitHub API 2026-07-08)
records, beyond a pull-request reference, only that it ``requires empirical
analysis''. The GGM depth $D$ and
arity $A$, which bound the number of epochs a note can ever traverse, are
symbolic throughout the book; the crate pins working values
(\S\ref{sec:tachyon-design-ggm}) that no design document yet ratifies. The ePrint note supplies only the
qualitative trade: with long epochs (``e.g., years'') the ecosystem ``can
perhaps organically support'' storage of past nullifiers, but payment-scale
throughput requires epochs ``perhaps even as short as a single block'',
which is precisely the regime that concentrates history into few providers
(ePrint 2025/2031, \S1.4).

\paragraph{What a specification must pin down.} The epoch length in blocks
and its activation semantics; $D$ and $A$, and therefore the lifetime
ceiling $A^D$ epochs a note can survive; the default and maximum delegation
window; and the reorganisation behaviour of the two-epoch consensus scan
(\S\ref{sec:tachyon-dbu-rollover}), which the book does not treat at all.

\paragraph{What would falsify the design.} The absence of a satisfying
point: if every epoch length either (short) grows the historical nullifier
database past what unincentivised parties store, or (long) grows the
per-note absence-proof work and time-to-spend past what wallets and services
sustain, then the parameter the design assumes exists does not. The
designers' own scaling narrative places the target regime at the short end,
where their own data-availability warning applies
(\S\ref{sec:tachyon-open-da}).

\subsubsection{The rollover soundness model}\label{sec:tachyon-dbu-rollover}

\paragraph{What exists.} A book-only consensus rule and a prose soundness
argument. The published design rejects note expiry --- the ePrint note
considers expiring notes and nullifiers and discards it because ``users must
move their funds into a new note, once per window, or lose their funds''
(ePrint 2025/2031, \S1.1) --- so any expiry-based description of Tachyon is
obsolete; notes stay spendable indefinitely via delegated absence proofs.
What replaces expiry is rollover: a spend publishes \emph{two} nullifiers,
for the anchor epoch and the next, and consensus checks every tachygram for
duplicates over exactly the current and immediately preceding epoch,
rejecting a duplicate within that two-epoch window
(\texttt{book/src/tachygrams.md} @ \texttt{26fdc165}). Everything older may
be discarded: ``Validators must only maintain the leading edge of revealed
nullifiers (to prevent duplicates in concurrent transactions) but can
otherwise permanently prune'' (forum t/50789, post \#22, 2025-11-04). The
ePrint note carries only the looser ``validators must still check for
duplicates''. The book's Soundness section argues the composition in prose:
splice identities (Construction~\ref{def:tachyon-design-splice})
checked at a Fiat--Shamir challenge, commit-equality binding of witnessed
sequences, and $\cmsym$-threading through the spendable lineage
(\texttt{book/src/proof-tree.md}). There are no formal definitions, no
theorems, and no machine-checked statements; the founding post itself
concedes that recursive-soundness bounds of PCD constructions are ``largely
ignored in practice'' (blog, 2025-04-02). The tracker carries live soundness
and malleability issues against the design: ``Bundle commitment is
order-free over actions, so bundles are malleable to action reordering''
(\#137), unchecked padding in proof encodings (\#161), and sub-block state
landing in stamp order rather than transaction order (\#135; all open,
fetched 2026-07-08). Even the rule's modality is unsettled: the
anchor chapter hedges, in an in-file comment, between ``may'', ``should'',
and ``must'' for which pool states consensus acknowledges
(\texttt{book/src/anchor.md}).

\paragraph{What a specification must pin down.} The consensus rule as
normative text: the exact scan set, the epoch-transition tolerance, whether
pruning is a permission or an obligation, and behaviour under
reorganisation across an epoch boundary; the tachygram-set commitment
format the scan runs over; and a soundness statement with a proof --- at
minimum, that two spends of one note accepted in any pair of epochs
necessarily collide inside some single two-epoch scan window, given the
dual-nullifier publication rule, and that the statement survives
composition with the recursive proof system's actual extraction bound.

\paragraph{What would falsify the design.} A double-spend across the
pruning boundary: two accepted spends of one note whose published nullifier
pairs never meet inside one scan window. The two-epoch window is the entire
replacement for the permanent nullifier set; a counterexample --- or a
demonstrated action-reordering exploit graduating from the open malleability
issues --- refutes the state-pruning claim that motivates the project.

\subsection{Open problems}\label{sec:tachyon-open}

Each entry below is acknowledged as unsolved by the designers; we quote the
acknowledgement. These are not gaps a specification could close by
transcription --- in each case the missing object is the construction
itself.

\subsubsection{Timing-correlation linkage}\label{sec:tachyon-open-timing}

\paragraph{What exists.} The problem statement, in the designers' words:
wallets hold multiple notes and update their unspent proofs together ---
``such as when first turned on after a period of inactivity, or charging on
WiFi at home. Transactions spending these notes to different parties could
still be linked together by their use of incremental unspent proofs that
were computed at correlated times. This is not a small problem; in fact, it
would be fatal to the entire approach'' (ePrint 2025/2031, \S1.3). The
remedy on offer is one clause --- re-randomizing the proof, by citation to
accumulation schemes (ePrint 2020/499) --- plus a gesture in the founding
post at ``network privacy countermeasures like mixnets'' and a promised
future post, which never appeared (blog index ends 2025-10-16).
\textsf{Ragu}'s \texttt{rerandomize} exists as code
(\S\ref{sec:tachyon-dbu-oss}), but no published argument connects it to the
property the ePrint needs: that a re-randomized proof is indistinguishable
from one the service never touched, under the timing side information the
service actually has. Collusion of a service with a network observer has no
published analysis anywhere in the corpus.

\paragraph{What a resolution must establish.} A formal indistinguishability
property for re-randomized PCD in \textsf{ragu}'s accumulation regime, a
proof that \texttt{rerandomize} achieves it, and --- separately, because
re-randomization does not touch timing --- a traffic discipline (batching,
cover queries, or a mixnet layer with a named threat model) that removes the
correlation channel the designers describe. The transport side is currently
one word in one blog post.

\paragraph{What would falsify the design.} The designers have already
stated it: a working correlation attack that links two spends through the
service that helped maintain both ``would be fatal to the entire
approach''. This is the only entry in the corpus where the falsification
criterion is supplied verbatim by the designers, and it remains open.

\subsubsection{Data availability, liveness, and incentives}
\label{sec:tachyon-open-da}

\paragraph{What exists.} A designer-acknowledged hazard with no mechanism.
Validators prune what wallets and services later need to prove absence over
(\S\ref{sec:tachyon-dbu-rollover}); the history must therefore live
somewhere off-consensus. The ePrint note describes the end state at scale:
``one could easily imagine a point where the historical nullifier database
grows so large that it can only be stored by a few highly centralized
providers \ldots\ the sudden shutdown of these services could render the
funds of all their dependent users unspendable, potentially permanently''
(ePrint 2025/2031, \S1.4). No new layer is planned: ``Tachyon will not
involve a new DA layer; the blockchain will continue to be used to post
nullifiers and note commitments, but not much else involving Tachyon''
(forum t/50789, post \#25, 2025-11-16); a DA layer ``would probably exist
solely to incentivize replication of the historic nullifiers/commitments,
and more importantly to incentivize running oblivious syncing services''
(ibid.). The designer's stated position is deferral: ``I'm looking forward
to talking about this more in the future, but it's not a concern of mine
right now'' (post \#22, 2025-11-04). The community objection is on the
record and unrebutted: building the non-inclusion proof ``requires
knowledge of every used nullifier'' (post \#19, 2025-10-29) --- someone must
keep everything consensus discards. No incentive design of any kind exists.

\paragraph{What a resolution must establish.} Who stores the pruned
nullifier and commitment history, under what incentive, with what
replication factor; a liveness bound --- how long a wallet's funds remain
spendable after its services vanish; and a recovery path for a user whose
absence-proof lineage lapses over a pruned interval. The ePrint's own
remedy is a sketch of a few paragraphs, framed as a possibility (``The
second possibility is to build a form of data-availability (DA) layer'',
\S1.4), not a design.

\paragraph{What would falsify the design.} The designers' scenario
realised: the throughput target forces short epochs
(\S\ref{sec:tachyon-dbu-epoch}), short epochs force a large historical
database, storage centralises, and a provider shutdown strands funds ---
``unspendable, potentially permanently''. A scaling design whose failure
mode at its own target scale is permanent loss of user funds, with no
incentive mechanism even sketched, is refuted by the first such event; the
open problem is to make the event impossible rather than unlikely.

\subsubsection{Recoverability}\label{sec:tachyon-open-recover}

\paragraph{What exists.} An acknowledged regression and a one-sentence
mitigation. Removing in-band secret distribution removes the property that
the chain itself is the backup: ``the user cannot rely on the blockchain as
a storage mechanism for recovering their funds from a seed phrase''
(blog, 2025-04-02). The designers accept the cost as forced: ``Getting this
right will be challenging, but we simply have no choice but to do it''
(forum t/50789, post \#22, 2025-11-04). Asked directly, the answer remains
non-committal: the replacement infrastructure ``opens up a can of worms
like who runs it? \ldots\ What if it's the people building it like don't
exist anymore? How do I get my funds?'' (ZK Podcast episode 388 transcript,
2026-01-21). The entire written resolution as of 2026-07-08 is the second
sentence of issue \#121: ``We're retaining in-band secret distribution for
recover from seed flows'' --- an issue whose first sentence declares the
book stale against it. The roadmap's alternative delivery path warns of
``backup/recoverability trade-offs''; a direct forum question about
seed-phrase recovery (2025-10-15) went unanswered.

\paragraph{What a resolution must establish.} A complete recover-from-seed
procedure: what the retained in-band channel carries and where it lives in
a protocol whose validators prune; how a restored wallet discovers its
notes' creation epochs and positions; and how it rebuilds an absence
lineage over history that consensus discarded --- which makes recoverability
a dependent of the data-availability problem
(\S\ref{sec:tachyon-open-da}), a dependency no document states.

\paragraph{What would falsify the design.} A seed-only restoration that
cannot terminate: a user with a correct seed phrase whose funds are
unrecoverable because the services holding the required history no longer
exist. For a system marketed as a payment protocol upgrade, demonstrated
non-recoverability at any realistic scale refutes the wallet model, not a
parameter of it.

\subsubsection{Tip-nullifier exposure versus the obliviousness claim}
\label{sec:tachyon-open-tip}

\paragraph{What exists.} A marketing claim and a design detail in tension.
The project overview states: ``the service never learns your actual
nullifiers, because the protocol forces them to periodically evolve in an
unlinkable way'' (tachyon.z.cash/overview, fetched 2026-07-08). The design
book's delegation construction hands the service one nullifier per
epoch-boundary crossing \emph{``plus the nullifier of the epoch in progress
at the span's tip, which is carried separately because that epoch is not
yet complete''} (\texttt{book/src/nullifiers.md} @ \texttt{26fdc165}). The
tip value is not a spent, superseded nullifier: the spendable's current
nullifier is ``the nullifier the wallet would publish to spend now''
(ibid.), so a service holding it can watch the pool and recognise the
client's spend in that epoch --- linking a service session to an on-chain
transaction, exactly the linkage obliviousness is sold as preventing. The
role table shows the escape: \textsf{UnspentSeed} and the other segment
steps are marked \emph{possible} for the wallet
(\texttt{book/src/proof-tree.md}), so a wallet can keep the leading edge
private and delegate only completed epochs. No document mandates that flow,
and none analyses the leak when it is not followed.

\paragraph{What a resolution must establish.} Either a mandated flow in
which the tip nullifier never leaves the wallet --- with the wallet-side
cost of sealing each epoch stated --- or a bound on what a service holding
the tip learns, under the timing side information of
\S\ref{sec:tachyon-open-timing}. Until one of these exists, the overview's
``never'' is falsified by the design book's own construction, and the
accurate statement is: the service never learns \emph{past} nullifiers,
and learns the current one unless the wallet does extra work no document
requires of it.

\paragraph{What would falsify the design.} Service-side spend detection in
the delegated flow --- trivially demonstrable today from the book's data
flow. The narrow claim is already false as advertised; what remains open is
whether the protocol will mandate the wallet-side variant that makes it
true.

\subsection{Quantitative claims without a published basis}
\label{sec:tachyon-figures}

The headline figures attached to Tachyon are not project measurements, and
a volume at this series' standard must say so plainly.

\begin{itemize}
\item The homepage states that Tachyon ``shrinks transactions by two orders
of magnitude and removes runaway state growth for validators''
(tachyon.z.cash, fetched 2026-07-08). No published benchmark, size
specification, or derivation supports the factor; it refers to marginal
on-chain size after aggregation and miner-side stripping, and no primary
document computes it.
\item The circulating throughput figures --- roughly $2{,}200$ TPS, or
$660$ TPS at ${\sim}1.5$ kB per transaction, and a $9.3$ kB${}\to{}450$ B
size reduction --- trace to Messari and CoinDesk analyst reports
(the latter commissioned by GenZcash, 2026-06-30), not to the project.
They are analyst estimates. Bowe's posts contain no TPS numbers, and the
founding post's promised PCD-toolkit benchmarks (2025-04-02) have not
appeared.
\item The only in-project performance numbers are \textsf{ragu}'s: verifier
time ``approximately 100ms'' at recursion threshold $K = 11$ and a
break-even against Orchard batch verification at roughly $50$ transactions
--- self-labelled ``rough estimates that should be validated with empirical
benchmarks'' (\textsf{ragu} book, \texttt{protocol/analysis.md} @
\texttt{830bbcda}). A criterion bench harness runs in CI and publishes no
numbers.
\item The claim that syncing services do only ``a similar amount of
computational effort (asymptotically)'' to today's validators is prefixed
``I believe'' in its only source (forum t/50789, post \#28, 2025-11-21);
no complexity analysis or benchmark exists.
\item The founding post's timeline --- ``many of the major scalability
improvements should be able to hit mainnet within a year'' (2025-04-02) ---
is falsified: as of 2026-07-08 nothing runs on testnet or mainnet, no ZIP
is upstreamed, the core ZIP stubs are empty, the crate's proof layer is
stubbed, and \textsf{ragu} is explicitly pre-audit. The current roadmap
carries no dates.
\end{itemize}

\subsection{Revisit triggers}\label{sec:tachyon-revisit}

The classifications above are dated 2026-07-08 and decay. Any of the
following moves material between bins and obsoletes parts of this section:
merge of ZIP-248 (the transaction-format substrate every Tachyon bundle
design defers to, open as \textsf{zcash/zips} PR \#1156); upstreaming of
the project's four fork amendment drafts (PRs \#1--\#4, open since
2026-06-02); a filled core ZIP stub (\#103--\#107); a published PIR payment
protocol; and \textsf{ragu}'s first published benchmark or audit. The
update-ZIP integrations remain undecided in the designers' own words ---
``still under discussion and nothing has been decided yet'' for the ZIP-221
history-tree leaf and the ZIP-317 fee contribution, ``blocked on ZIP-248
reaching editor consensus'' for bundle registration
(\texttt{book/src/zips/} @ \texttt{26fdc165}).

\section{Security posture and outlook}\label{sec:tachyon-outlook}

This closing section states the volume's synthesis. We compress the
preceding chapters into a single conservation law
(\S\ref{sec:tachyon-conservation}), place Tachyon at its one point of
contact with the post-quantum programme
(\S\ref{sec:tachyon-pq-intersection}), fix its relation to the surrounding
consensus roadmap (\S\ref{sec:tachyon-environment}), and end with the
volume's verdict: the concrete artifacts that must exist before any Tachyon
claim can be checked rather than believed
(\S\ref{sec:tachyon-checkability}). Ground truth is pinned as throughout the
volume: the design book at \texttt{tachyon-zcash/tachyon} @
\texttt{26fdc165} (2026-07-07), Ragu at \texttt{tachyon-zcash/ragu} @
\texttt{830bbcda} (2026-07-05), the ZIP corpus at \texttt{zcash/zips} @
\texttt{c6b70358} (2026-07-05); web sources carry individual dates. The
material compressed here is constructed in \S\ref{sec:tachyon-fold} (the
proof-carrying fold and its nineteen steps), \S\ref{sec:tachyon-dbu-nf}
(nullifier evolution), and \S\ref{sec:tachyon-dbu-oss} (the sync-service
model); Ragu's code state is \S\ref{sec:tachyon-ragu}; the open problems are
\S\ref{sec:tachyon-open}.

\subsection{The conservation law}\label{sec:tachyon-conservation}

The designers claim the scaling gain costs nothing in the two properties
Zcash exists to provide. The eprint promises ``continual, permanent pruning
of nullifiers by validators'', achieved ``without compromising privacy'';
its syncing services are \emph{untrusted}, hold only ``ephemeral state per
client'', and ``cannot link their clients to any transactions that
ultimately appear on the public ledger'' (Bowe--Miers, eprint 2025/2031,
abstract; revised 2025-11-03). Nothing in the design asks a user to trust a
third party with funds or with integrity --- a service that lies is caught
by proof verification --- and the assumption family, Pallas/Vesta discrete
logarithms plus random-oracle heuristics, is the one already deployed in
Orchard. In this exact sense soundness and privacy are preserved \emph{in
principle}.

The qualification carries weight. Soundness of the proof tree is argued in
prose (the book's Soundness section, \texttt{book/src/proof-tree.md} @
\texttt{26fdc165}; the eprint contains zero theorems and no formal security
definitions); the term ``oblivious'' has no security definition anywhere in
the corpus; the hiding of the new Poseidon-based commitments is asserted,
not proven; and the founding post itself concedes that the recursive
soundness bounds of PCD constructions are ``largely ignored in practice''
(Bowe, ``Tachyon: Scaling Zcash with Oblivious Synchronization'',
2025-04-02). Each of these is \emph{designed but unspecified} at best, and
none changes bins in this section.

What the design spends is availability. Pruning does not destroy the
history a wallet needs; it relocates the obligation to keep it. Validators
drop exactly what wallets later require to prove their notes unspent, and
someone must hold it. The designers rule out a protocol-level home for that
obligation: ``Tachyon will not involve a new DA layer; the blockchain will
continue to be used to post nullifiers and note commitments, but not much
else involving Tachyon'' --- while observing that under Tachyon users
``become more sensitive to the availability of the history in order to make
their funds spendable'', and that a data-availability layer, if one arose,
``would probably exist solely to incentivize replication'' (forum t/50789
\#25, 2025-11-16). The computation relocates the same way: the expectation
that syncing services will spend ``a similar amount of computational effort
(asymptotically) that validators currently do'' is a designer's ``I
believe'' (forum t/50789 \#28, 2025-11-21) with no published complexity
analysis behind it.

We compress the relocation into a pick-two. Deployed Orchard provides three
availability properties simultaneously --- and pays for them with the
unbounded validator state Tachyon exists to remove:
\begin{enumerate}[label=(\roman*)]
\item \emph{full obliviousness}: no helper party learns anything about the
  wallet's notes or transactions;
\item \emph{no liveness dependence}: spendability does not require any
  particular external party to be running;
\item \emph{chain as complete backup}: a seed phrase plus the public chain
  recover all funds.
\end{enumerate}
Once validators prune, a wallet retains at most two. Keeping (ii) and (iii)
means the chain retains everything a recovering wallet needs --- that is
today's protocol, not Tachyon. Keeping (i) and (ii) means the wallet
ingests and stores the pruned nullifier history itself, continuously; the
chain stops being its backup, exactly as the founding post concedes: ``the
user cannot rely on the blockchain as a storage mechanism for recovering
their funds from a seed phrase'' (2025-04-02). Keeping (iii) through
services --- delegated absence proofs reconstructing spendability from seed
--- surrenders (ii): ``the sudden shutdown of these services could render
the funds of all their dependent users unspendable, potentially
permanently'' (eprint 2025/2031, \S1.4). And the retained (i) is itself the
design's most fragile member: correlated proof-update timing linking a
wallet's transactions ``is not a small problem; in fact, it would be fatal
to the entire approach'', with the remedy --- proof re-randomization ---
supplied by citation rather than construction (eprint 2025/2031, \S1.3);
this volume classifies it an \emph{open problem}.

The project's own trajectory moves inside this triangle rather than out of
it. The May-2026 pivot retains ``in-band secret distribution for recover
from seed flows'' (tachyon-zcash/tachyon issue \#121, 2026-05-26), buying
back the note-discovery half of (iii) at the price of on-chain ciphertexts
for those flows, while spendability from seed still rests on delegated
absence proofs and therefore on (ii); the roadmap's fallback --- dedicated
note-delivery infrastructure --- concedes ``UX or backup/recoverability
trade-offs'' in as many words (tachyon.z.cash/roadmap, fetched 2026-07-08).
Classification of the trilemma itself: the framing is this volume's; every
fact inside it is designer-acknowledged and individually cited; and no
Tachyon document states the trade-off in one place. A specification will
have to, because the chosen corner determines what a wallet must store ---
the difference between a seed phrase that recovers funds and one that does
not.

\subsection{The post-quantum intersection}
\label{sec:tachyon-pq-intersection}

The post-quantum analysis of this series (the \emph{Ironwood Guide},
\S``Post-quantum security of the Orchard protocol'') reaches one headline
conclusion about deployed Orchard: exactly one quantum
threat is retroactive. Note encryption is the canonical
harvest-now-decrypt-later target --- every
$(\mathsf{epk}, C_{\mathrm{enc}})$ pair recorded on chain today becomes
decryptable by a future discrete-log-capable adversary who also holds the
recipient's address, exposing amounts, memos, and the full receiving
history, permanently --- while every integrity break (Sinsemilla binding,
value-commitment binding, RedPallas forgery, Halo~2 soundness) is
prospective and mitigable by switching pools off in time. The deployed
mitigation programme declines to touch the retroactive item: ``The
situation with respect to Privacy is unchanged for any pool''; the note
encryption of the deployed pools ``is subject to `Harvest Now, Decrypt
Later' attacks'', with countermeasures only ``under consideration'' (ZIP
2005, Non-requirements and Privacy discussion, zips @ \texttt{c6b70358}).

Tachyon is the only artifact in the corpus that addresses this liability,
and it does so structurally rather than cryptographically: it removes the
ciphertext. ``We plan to remove in-band secret distribution from Tachyon
transactions to fully defend against all quantum attacks on privacy in
on-chain transactions. This step is essential, because harvest now, decrypt
later threats are perilous for blockchain protocols'' (Bowe, ``Zcash and
Quantum Computers'', 2025-10-16). Where no key-agreement ciphertext is
published, there is nothing to harvest, and the remaining on-chain objects
--- hiding commitments, zero-knowledge proofs, re-randomized keys --- leak
nothing to a quantum adversary who lacks the address. Three qualifications
bound the claim. First, the defence is prospective only: ciphertexts
already recorded remain harvestable forever; no protocol change can repair
them. Second, the pivot of issue \#121 retains in-band distribution for
recover-from-seed flows, and whether those ciphertexts use the deployed
discrete-log key agreement or the ``post-quantum key exchange'' the roadmap
names for its PIR payment protocol is pinned down nowhere --- the PIR
protocol has no published specification at all (tachyon.z.cash/roadmap,
fetched 2026-07-08). Third, no post-quantum privacy analysis of Tachyon
exists; the statistical hiding of its Poseidon-based commitments is
asserted, not proven. Bin: \emph{designed but unspecified}, and the
recovery-flow cryptography is the first thing a specification must fix
before the post-quantum privacy claim becomes evaluable.

On the soundness side the same project deepens the discrete-logarithm
commitment. Ragu is an ECDLP-based recursive SNARK by declared design: the
requirement of small recursive proofs ``essentially limits us to
ECDLP-based constructions with our existing payment protocol
architecture'', with a footnote recording ``reduced concern (in the medium
term) for post-quantum \emph{soundness} risks'' that cites the same blog
post (\texttt{book/src/protocol/index.md}, ragu @ \texttt{830bbcda}). The
designers' asymmetry argument is explicit: ``Quantum privacy breaks are
usually retroactive, whereas soundness breaks are usually not'' (Bowe,
2025-10-16); a prospective soundness failure is answerable by reaction ---
turnstiles bound inflation, and ZIP 2005 re-anchors note ownership in
symmetric primitives so that funds survive a migration. The argument is
coherent, and it prices one thing incompletely: depth. Under Tachyon the
discrete-log assumption no longer underwrites a single Action circuit but
the recursion substrate itself, and the wallet's state --- the
proof-carrying data that replaces the chain as the user's record --- is a
stack of DL-sound proofs whose recursive soundness bounds the founding post
already concedes are ``largely ignored in practice'' (2025-04-02). A future
migration off the assumption is then a migration of the state model, not
the replacement of one circuit. Bin: \emph{open problem}, jointly with the
unquantified recursion-soundness bound; no document in the corpus analyses
either.

\subsection{Relation to Crosslink and the upgrade environment}
\label{sec:tachyon-environment}

Crosslink --- the finality/proof-of-stake effort led by Shielded Labs ---
and Tachyon are formally independent. The founding post sees ``no reason to
believe that Tachyon will conflict with any of the active areas of
development such as Crosslink and ZSAs'' (2025-04-02);
Shielded Labs states that ``Crosslink and Tachyon are separate efforts that
don't necessarily interact with each other with regards to roadmapping +
timelines'' (forum t/50789 \#11, 2025-10-28); and no interaction analysis
exists in either direction --- a repository-wide search of the ZIP corpus
finds Crosslink only as a citation footnote in ZIP 218 (25-second block
target spacing, Status: Draft, created 2026-03-13; zips @
\texttt{c6b70358}, searched 2026-07-08). The independence is therefore
asserted, not examined; both proposals alter the environment the other
assumes (consensus finality and block cadence on one side, validator state
and epoch-windowed double-spend checking on the other), and Tachyon's epoch
period is itself untuned (tachyon-zcash/tachyon issue \#79). The near-term
environment is fixed without Tachyon: NU6.2 is settled (ZIP 257, Status:
Final --- the emergency Action-circuit correction), NU6.3 carries the
Ironwood pool and the ZIP 2005 recoverability programme (ZIPs 258, 229,
2005; Draft/Draft/Proposed), and Tachyon's own standing is a mandate
without a vehicle: ``universal or near-universal support across every group
\ldots{} the most unambiguous signal from the entire poll'' in the February
2026 NU7 sentiment poll, with the stated next step only to ``assess
technical readiness and scope for inclusion in NU7'' (zfnd.org, NU7 Polling
Results, 2026-02-23), while the deployment draft
(\texttt{draft-arya-deploy-nu7.md}, heights TBD, zips @ \texttt{c6b70358})
lists no Tachyon ZIP. Beyond NU7 there is nothing to relate Tachyon to: the
ZIP corpus contains no NU8 artifact of any kind (repository-wide search of
\texttt{zcash/zips} @ \texttt{c6b70358}, 2026-07-08); the relation to NU8
becomes statable only when an NU8 deployment draft exists.

\subsection{The checkability boundary}\label{sec:tachyon-checkability}

Every claim in this volume carries one of three labels: \emph{specified},
\emph{designed but unspecified}, or \emph{open problem}. The distribution
of those labels is the verdict. The specified bin contains standing and
environment facts only --- the poll result, editor affiliations, the
NU6.2/NU6.3 context, the ZIP 321/324 out-of-band substrate, and Ragu's code
state --- and not one byte format, consensus rule, or security theorem of
Tachyon itself. Three artifact classes would move the boundary.

\paragraph{Benchmarks.} No performance figure exists at project rigor. The
founding post promised benchmarks of a proof-carrying-data toolkit ``to set
a floor on the performance of shielded transaction aggregation and
oblivious syncing services'' (2025-04-02); none has appeared. The only
in-project numbers --- roughly $100$~ms aggregated verification at
recursion threshold $K = 11$ and a break-even against Orchard batch
verification near $50$ transactions --- are self-labelled ``rough estimates
that should be validated with empirical benchmarks''
(\texttt{book/src/protocol/analysis.md}, ragu @ \texttt{830bbcda}); the CI
bench harness publishes no output; the circulating throughput and size
figures are analyst estimates, not project figures; and the homepage's
two-orders-of-magnitude headline has no published derivation. A meaningful
benchmark moreover requires code that does not yet exist: Ragu's
\texttt{Proof} carries full polynomials, verification is linear-time, and
no inner-product opening argument or decider is implemented
(\texttt{crates/ragu\_pcd/src/\{proof/mod.rs,verify.rs\}} @
\texttt{830bbcda}). Until a benchmark of non-mock Ragu PCD with succinct
proofs exists, every scaling claim is unfalsifiable.

\paragraph{The nullifier-evolution specification.} The one genuinely novel
cryptographic object --- the evolving nullifier --- exists at
eprint-sketch and design-book rigor only: the eprint offers a PRF sketch
``as a concrete example'' (2025/2031, \S1.3); the concrete GGM/Poseidon
derivation, the dual-nullifier spend, and the two-epoch duplicate window
live only in the book (\texttt{book/src/\{nullifiers,tachygrams\}.md} @
\texttt{26fdc165}); the crate stubs the derivation and calls the commitment
scheme ``TBD'' (\texttt{crates/tachyon/src/note.rs}); the derivation is
churning through open alternatives (issues \#139, \#86, \#87/\#97/\#98)
with live soundness and malleability issues (\#161/\#137/\#135); and the
epoch period, GGM depth, and arity on which the entire pruning model hinges
are untuned (issue \#79, ``requires empirical analysis''). A specification
must pin the derivation formulas and domain separators, the epoch
parameters, the dual-nullifier consensus rule and the pruning rule, and ---
above all --- the security definitions (per-epoch unlinkability, service
obliviousness) that nothing in the corpus states formally.

\paragraph{Core ZIPs.} The consensus path is empty. The five core ZIPs that
would carry the protocol --- Shielded Protocol (\#103), Bundle (\#104),
Accumulator (\#105), Aggregator (\#106), OSS (\#107) --- are headers-only
stubs of $216$--$245$ bytes each (\texttt{book/src/zips/} @
\texttt{26fdc165}); the four peripheral amendment drafts sit unmerged on
the project's own fork (tachyon-zcash/zips PRs \#1--\#4, opened 2026-06-02,
open as of 2026-07-08); upstream \texttt{zcash/zips} carries zero Tachyon
ZIPs; and the extensible-format substrate everything defers to, ZIP 248,
is itself an unmerged PR idle since 2026-05-26 (zcash/zips\#1156). Beside
the ZIPs stand Ragu's own declared gates: a Fiat--Shamir domain tag that is
literally \texttt{b"FIXME"} (\texttt{crates/ragu\_pcd/src/lib.rs:50--51}),
registry binding flagged \texttt{FIXME(security)} as insufficient
(\texttt{crates/ragu\_circuits/src/registry.rs:673--692}), three missing
verifier checks (\texttt{crates/ragu\_pcd/src/verify.rs:122}), and no
audit: ``Ragu is under heavy development and has not undergone auditing''
(\texttt{README.md} @ \texttt{830bbcda}).

The events that would obsolete this volume's classification are the revisit
triggers of \S\ref{sec:tachyon-revisit}, read forward. Until they occur, the
posture admits a one-sentence summary:
soundness and privacy are conserved in intention, availability is spent by
design, and none of it is yet checkable.

\end{document}
